\chapter[Phase 10 closures]{Phase 10: closures, reductions and examples (other results)}\label{ch:cIV4}
\chapterbrief
{the rows of the ledger sections F10-A, F10-B2, F10-C2, F10-COL, F10-GTU, F10-INT, F10-LEAN and F10-TRIAGE. Where two sections give the same identifier, the later one prevails (F10-B2 over F10-B, F10-C2 over F10-C, F10-TRIAGE over all).}
{to prove, in full, the Phase~10 results whose argument is self-contained and short enough to write out (the dilation theorem, the closedness lemma, the annulus counterexample, the Dyson--Phillips regularity, the rank-$2$ counterexample and the Cantor-type zeros), and to record the others with the status exactly as in the ledger.}
{Dim.~B117--B123, then list cards. The two proofs that change statements of Phases~0--5 are A52 and A53 (Chapter~\ref{ch:cIII6}). The summary of the results is Chapter~\ref{PAP-ch:phase10} of the paper.}

\section{Non-commuting flows and the factor decomposition}

\begin{dimcard}{B117}{S10A-T2, P1-Q7}{No finite word of flows at fixed coupling, under a weighted dilation}{\pw}{ch:phase10}
Let $\mathcal A$ be a unital associative $\C$-algebra, $X,Y\in\mathcal A$ satisfying hypothesis (H) of P1-T5 (the elements $\mathrm{ad}_X^kY$ are linearly independent and $X\ne0$), and $\mathcal A_0$ the unital subalgebra they generate. Assume \textbf{(Dil)}: there are integers $\alpha\ge0$, $\beta<\alpha$ and an infinite $\Lambda\subset\C^*$ such that for each $\mu\in\Lambda$ a unital endomorphism $\sigma_\mu$ of $\mathcal A_0$ has $\sigma_\mu X=\mu^\alpha X$, $\sigma_\mu Y=\mu^\beta Y$. Then for every $\varepsilon_0\in\C^*$ there are no word $(Z_1,\dots,Z_d)\in\{X,Y\}^d$ and series $g_i\in t\C[[t]]$ with $e^{t(X+\varepsilon_0Y)}=\prod_{i=1}^de^{g_i(t)Z_i}$ in $\mathcal A[[t]]$. Hence P1-T5 holds without the hypothesis ``the flows of $Y$ vanish at $\varepsilon=0$''.
\end{dimcard}

\begin{proofidea}
Rescale by $\sigma_\mu$ and substitute $t=\mu^{-\alpha}s$: the fixed coupling $\varepsilon_0$ becomes $\varepsilon_0\nu^{\mathfrak a}\to0$ and the flows of $Y$ become $1+O(\nu^{\mathfrak a})$. The old proof of P1-T5 then works in $\C[\nu]$.
\end{proofidea}

\begin{proof}
\step{1}{Rescaling}
Applying $\sigma_\mu$ coefficientwise gives $e^{t(\mu^\alpha X+\varepsilon_0\mu^\beta Y)}=\prod_ie^{g_i(t)\mu^{w_i}Z_i}$ with $w_i=\alpha$ for $Z_i=X$ and $\beta$ for $Z_i=Y$. Put $t=\nu^\alpha s$, $\nu=\mu^{-1}$ and $\mathfrak a=\alpha-\beta\ge1$:
\[
e^{s(X+\varepsilon_0\nu^{\mathfrak a}Y)}=\prod_ie^{h_i(s,\nu)Z_i},\qquad h_i=\sum_kg_{i,k}\nu^{\alpha k-w_i}s^k.
\]
For $Z_i=X$ the exponents are $\alpha(k-1)\ge0$; for $Z_i=Y$ they are $\ge\alpha-\beta=\mathfrak a$.

\step{2}{A polynomial identity in $\nu$}
The coefficient of $s^m$ lies in the finite-dimensional span of words of length $\le m$ and has coordinates polynomial in $\nu$ on both sides. They agree on the infinite set $\Lambda^{-1}$, hence as polynomials.

\step{3}{Orders below $\nu^{\mathfrak a}$}
Modulo $\nu^{\mathfrak a}$ the $Y$-factors are the identity and the left side is $e^{sX}$; the $X$-factors commute, so $e^{sX}\equiv e^{HX}$ with $H=\sum_{Z_i=X}h_i$. If $H-s\not\equiv0$, its lowest term $c\,s^p\nu^q$ gives $cX=0$ in $e^{(H-s)X}=1$, impossible. So $H\equiv s$ mod $\nu^{\mathfrak a}$.

\step{4}{Order $\nu^{\mathfrak a}$}
On the left the coefficient is $\varepsilon_0e^{sX}\sum_k(-1)^k\frac{s^{k+1}}{(k+1)!}Y_k$ (formal Duhamel), $Y_k=\mathrm{ad}_X^kY$. On the right only two terms contribute: all $Y$-factors at order $0$ (giving $r(s)Xe^{sX}$, $r$ the $\nu^{\mathfrak a}$-coefficient of $H$), or exactly one $Y$-factor at order $\nu^{\mathfrak a}$ with coefficient $q_i(s)Y$. Multiplying by $e^{-sX}$ and using $e^{-\tau X}Ye^{\tau X}=\sum(-1)^k\frac{\tau^k}{k!}Y_k$:
\[
rX+\sum_k(-1)^kY_k\sum_{i:Z_i=Y}q_i\frac{\tau_i^k}{k!}=\varepsilon_0\sum_k(-1)^k\frac{s^{k+1}}{(k+1)!}Y_k.
\]

\step{5}{The Hankel contradiction}
By (H), $r=0$ and $\sum_{i:Z_i=Y}q_i\tau_i^k=\varepsilon_0s^{k+1}/(k+1)$ for all $k\ge0$. With $p=\#\{i:Z_i=Y\}$ and $n=p+1$, the Hankel matrix $[m_{i+j}]_{i,j<n}$ of the left sides is $W^{\mathsf T}\mathrm{diag}(q)W$ with $W$ of size $p\times n$, so its determinant vanishes identically (Cauchy--Binet). That of the right sides is $\varepsilon_0^ns^{n^2}\det\mathrm{Hilb}_n\ne0$. For $p=0$, $m_0=\varepsilon_0s\ne0$. Contradiction.

\whyworks{dilation turns a fixed coupling into one that tends to $0$, recovering by symmetry the hypothesis that was to be removed.}
\end{proof}
\uses{A24.}

\begin{caution}[title=What remains open]
Algebras with (H) but without weighted dilations. The candidate $X=(x^3+x)\partial$, $Y=\partial$ in the Weyl algebra $A_1$ is not homogeneous; it was settled in the triage, but the general case stays open (P1-Q7).
\end{caution}

\begin{dimcard}{B118}{S10A-L2, C2}{A factor with compact interaction closes the sum}{\pw}{ch:phase10}
Let $\mathcal H_1,\dots,\mathcal H_r$ be closed subspaces of a Hilbert space with orthogonal projections $\Pi_i$. If $M=\mathcal H_1+\dots+\mathcal H_{r-1}$ is closed and $\Pi_r\Pi_i$ is compact for each $i<r$, then $\mathcal H_1+\dots+\mathcal H_r$ is closed. For $r=3$, put an edge $\{i,j\}$ when $\Pi_i\Pi_j$ is not compact: if the three pairwise sums are closed and the triple sum is not, every vertex has degree $\ge1$, so there are at least two edges.
\end{dimcard}

\begin{proof}
\step{1}{The lemma}
The sum map $S:\bigoplus_{i<r}\mathcal H_i\to M$ is bounded and onto the Banach space $M$. By the open mapping theorem each $m\in M$ with $\|m\|\le1$ is $\sum_{i<r}h_i$ with $\sum\|h_i\|\le C$. Then $\Pi_rm=\sum_i\Pi_r\Pi_ih_i$, so $\Pi_r(B_M)\subset\sum_{i<r}\Pi_r\Pi_i(B_C)$, relatively compact. Hence $\Pi_r|_M$ and $\Pi_r\Pi_M$ are compact, and by OPR2-L1 for two closed subspaces $\mathcal H_r$ and $M$ the sum $\mathcal H_r+M$ is closed.

\step{2}{The corollary}
The relation is symmetric, $(\Pi_i\Pi_j)^*=\Pi_j\Pi_i$. If a vertex $k$ were isolated, apply the lemma with $\mathcal H_r=\mathcal H_k$ and $M=\mathcal H_i+\mathcal H_j$ (closed by hypothesis): the triple sum is closed, a contradiction. Three vertices of degree $\ge1$ need at least two edges.

\whyworks{compactness of the interaction lets the new subspace be handled as a finite-rank perturbation of an orthogonal complement.}
\end{proof}
\uses{none.}

\begin{caution}[title=Consequence for C12-Q2]
A counterexample to ``pairwise closed implies triple closed'' for $P=Q_1Q_2Q_3$ must have at least two pairs of factors with non-compact interaction. Real transports are excluded (OPR2-T2), and no candidate family has been analysed: the question stays open, reduced to such triples.
\end{caution}

\begin{dimcard}{B119}{S10A-P1-annulus}{The annulus: closed but not dense}{\pw}{ch:phase10}
Let $P=\Delta$, factors $Q_1=\partial_t+i\partial_x$, $Q_2=\partial_t-i\partial_x$ ($\delta(P)=1$), so $\mathcal H_1$ are holomorphic and $\mathcal H_2$ anti-holomorphic functions in $H^m$. (i) On the annulus $A=\{1<|z|<2\}$ (Lipschitz, non-convex, $b_1=1$), $\log|z|\in\ker_{H^m(A)}\Delta\setminus\overline{\mathcal H_1+\mathcal H_2}$ for every $m\ge0$, and for $m\ge1$ $\ker_{H^m(A)}\Delta=(\mathcal H_1+\mathcal H_2)\oplus\C\log|z|$: the sum is closed of codimension $1=b_1$ and not dense. (ii) If a bounded Lipschitz $\Omega$ has a cycle $\gamma$ and a point $a\notin\Omega$ with non-zero winding number, then $\log|z-a|\notin\overline{\mathcal H_1+\mathcal H_2}$. So P2-Q1 in its universal form is false, and the differential example ``closed but not dense'' of C12-P7 exists.
\end{dimcard}

\begin{proof}
\step{1}{The flux functional}
For a smooth cycle $\gamma\Subset\Omega$ and $h$ harmonic put $\Phi_\gamma(h)=\int_\gamma{*dh}$. With $\delta=\mathrm{dist}(\gamma,\partial\Omega)/2$, the mean-value property and Cauchy's gradient estimate give $|\nabla h(y)|\le C\delta^{-2}\|h\|_{L^2(\Omega)}$ on $\gamma$, so $|\Phi_\gamma(h)|\le C_\gamma\|h\|_{H^m}$: $\Phi_\gamma$ is continuous.

\step{2}{It vanishes on $\mathcal H_1+\mathcal H_2$}
For $f$ holomorphic, ${*df}=-if'dz$ and $\Phi_\gamma(f)=-i\oint_\gamma f'dz=0$ since $f$ is a single-valued primitive; for $\bar g$ anti-holomorphic, $\Phi_\gamma(\bar g)=\overline{\Phi_\gamma(g)}=0$. By continuity $\Phi_\gamma$ vanishes on the closure.

\step{3}{It does not vanish on the logarithm}
${*d\log|z-a|}=d\arg(z-a)$, so $\Phi_\gamma(\log|z-a|)=2\pi\,\mathrm{Ind}_\gamma(a)\ne0$. For a hole one takes $a$ in the interior of the bounded component of $\C\setminus\Omega$ (the index is constant on it); such a cycle exists when $b_1\ge1$ (Rudin RCA 13.11, \cf). This proves (ii) and the first part of (i) with $\gamma=\{|z|=\frac32\}$.

\step{4}{Exact codimension on the annulus, $m\ge1$}
A harmonic $h$ on $A$ is $f+\bar g+c\log|z|$ with $f,g$ holomorphic single-valued (Laurent) and $c=\Phi(h)/2\pi$. If $\Phi(h)=0$, then $h=f+\bar g$ and $f'=\partial_zh\in H^{m-1}$. In polar coordinates $|f(re^{i\theta})-f(\frac32e^{i\theta})|^2\le\int_1^2|f'|^2\varrho\,d\varrho$, so $\|f-f(\frac32e^{i\cdot})\|_{L^2(A)}\le C\|f'\|_{L^2}$, while $f(\frac32e^{i\theta})$ is bounded. Derivatives of order $1,\dots,m$ of $f$ are derivatives of $f'\in H^{m-1}$, so $f\in H^m$; likewise $\bar g$. Thus $\ker\Phi=\mathcal H_1+\mathcal H_2$, closed as the kernel of a continuous functional.

\whyworks{a hole carries a flux that no holomorphic or anti-holomorphic function can have, and that flux is a continuous functional.}
\end{proof}
\uses{A52, A53.}

\begin{dimcard}{B120}{S10A-T3}{Density up to the first Betti number, on Lipschitz regions}{\modulo{[F1], [F2]}}{ch:phase10}
Let $\Omega\subset\R^2$ be bounded, connected and Lipschitz, $P=c\prod_{j=1}^NQ_j$ with pairwise coprime primary $Q_j$, $\mathcal H_j=\ker_{H^m(\Omega)}Q_j$, $\mathcal K=\ker_{H^m(\Omega)}P$, $S=\overline{\sum\mathcal H_j}$. Then (a) $\mathrm{codim}_{\mathcal K}S\le\delta(P)\,b_1(\Omega)$; (b) if $\Omega$ is simply connected, $\sum\mathcal H_j$ is dense in $\mathcal K$, and the consequences of OPF5B-T1 hold.
\end{dimcard}

\begin{proofidea}
Approximate a solution in $\Omega$ by one on a neighbourhood $U\supset\bar\Omega$ with the same topology; there the cokernel of the sum map is bounded by $\delta(P)b_1$ (OPF-T4); a finite-dimensional argument transfers the bound back to $\Omega$.
\end{proofidea}

\begin{proof}
\step{1}{A uniform segment condition replaces convexity}
Cover $\partial\Omega$ by finitely many cylinders $C_i$ in which $\Omega$ is the supergraph of an $L_i$-Lipschitz function in direction $e_i$; choose $C_i'\Subset C_i$ still covering, a partition of unity $\theta_0,\theta_i$ and translations $\tau_i=\varepsilon e_i$. For $y_0\in\bar\Omega\cap\bar C_i'$ one has $B(y_0+\tau_i,r_\varepsilon)\subset\Omega$ with $r_\varepsilon=\varepsilon/(2\sqrt{1+L^2})$, because $y_2-\varphi_i(y_1)\ge\varepsilon-\sqrt{1+L^2}|q|>\varepsilon/2$. This is the only place where OPF5B-T1 used convexity.

\step{2}{Approximation by solutions on a neighbourhood}
The remaining steps of OPF5B-T1 (Stein extension [F2], continuity of translations, the facts F1 and L1, mollification) do not use convexity. For $u\in\mathcal K$ and $\epsilon>0$ there are an open $V'_\varepsilon\supset\bar\Omega$ and $v_\eta\in C^\infty$ with $P(\partial)v_\eta=0$ on $V'_\varepsilon$ and $\|v_\eta-u\|_{H^m(\Omega)}<2\epsilon$.

\step{3}{A neighbourhood with the same topology}
By the collar fact [COL] there is a flow $\Phi_\sigma$ of a transversal field with $\bar\Omega\subset\Phi_\sigma(\Omega)\subset\{\mathrm{dist}<C\sigma\}$; take $\sigma$ with $C\sigma<r_\varepsilon/4$ and $U=\Phi_\sigma(\Omega)$, diffeomorphic to $\Omega$, so $b_1(U)=b_1(\Omega)$.

\step{4}{The cokernel on $U$}
OPF-T4(ii) for $C^\infty$ gives $\dim\mathrm{coker}\,\Sigma_U\le\delta(P)b_1(U)=:\Delta$ with $\Sigma_U:\bigoplus\ker_{C^\infty(U)}Q_j\to\ker_{C^\infty(U)}P$. If $\Delta=0$, $v_\eta|_U=\sum w_j$ and, $\bar\Omega\subset U$ being compact, $w_j|_\Omega\in\mathcal H_j$: this is (b).

\step{5}{Finite codimension}
Suppose $\mathrm{codim}_{\mathcal K}S>\Delta$: there are $u_0,\dots,u_\Delta\in\mathcal K$ with independent images in the Hilbert space $\mathcal K/S$. Independence is open (Gram determinant), so there is $\epsilon$ with the same property for every family $\epsilon$-close. Steps 2--3 with one $U$ give $v_k\in\ker_{C^\infty(U)}P$ with $\|v_k|_\Omega-u_k\|<\epsilon$. As $\dim\mathrm{coker}\,\Sigma_U\le\Delta$, the $\Delta+1$ classes of $v_k$ are dependent: $\sum c_kv_k=\sum w_j$ on $U$ with $c\ne0$. Restricting to $\Omega$, $\sum c_kv_k|_\Omega\in\sum\mathcal H_j\subset S$, so the images in $\mathcal K/S$ are dependent. Contradiction.

\whyworks{a Lipschitz boundary supplies a uniform inward segment, a collar supplies a neighbourhood with the same holes, and the cokernel of the sum map there is bounded by the number of holes times $\delta(P)$.}
\end{proof}
\uses{A53.}

\section{Variable coefficients: regularity, counterexample and zeros}

\begin{dimcard}{B121}{VAR-P11(e)}{Classical solutions for a $C^1$ perturbation}{\pw}{ch:varcoef}
Let $A_0$ generate a $C_0$-semigroup $T(t)$ on a Banach space $X$, $\|T(t)\|\le M_0$ on $[0,T]$, $B\in C^1([0,T];\mathcal L(X))$ with $\beta_0=\sup\|B\|$, $\beta_1=\sup\|B'\|$, $U(t,s)$ the Dyson--Phillips propagator (the mild solution $u(t)=T(t-s)x+\int_s^tT(t-r)B(r)u(r)dr$). If $x\in D(A_0)$ then $u=U(\cdot,s)x\in C^1([s,T];X)\cap C([s,T];D(A_0))$ and $u'=(A_0+B(t))u$; the solution is unique.
\end{dimcard}

\begin{proof}
\step{1}{Convolution lemma}
If $f\in C^1$ and $F(t)=\int_s^tT(t-r)f(r)dr$, then $F\in C^1$, $F'(t)=T(t-s)f(s)+\int_s^tT(t-r)f'(r)dr$, and $F(t)\in D(A_0)$ with $A_0F=F'-f$. Write $F(t)=\int_0^{t-s}T(\sigma)f(t-\sigma)d\sigma$; the difference quotient has a boundary term tending to $T(t-s)f(s)$ and an integral term tending to $\int T(\sigma)f'(t-\sigma)d\sigma$ by dominated convergence. For the domain, $h^{-1}(T(h)-I)F(t)=h^{-1}[F(t+h)-F(t)]-h^{-1}\int_t^{t+h}T(t+h-r)f(r)dr\to F'(t)-f(t)$.

\step{2}{A candidate derivative}
With $y_0=A_0x+B(s)x$ let $v$ be the unique solution of $v(t)=T(t-s)y_0+\int_s^tT(t-r)[B'(r)u(r)+B(r)v(r)]dr$ (Picard, kernel bounded by $M_0\beta_0$). Put $w(t)=x+\int_s^tv$, $f=Bw$ (so $f'=B'w+Bv$) and $z(t)=T(t-s)x+\int_s^tT(t-r)f(r)dr$.

\step{3}{Gronwall}
By Step~1, $z'=T(t-s)(A_0x+B(s)x)+\int_s^tT(t-r)[B'w+Bv]dr$, so with $e=u-w$: $u-z=\int_s^tT(t-r)Be\,dr$ and $(w-z)'=v-z'=\int_s^tT(t-r)B'e\,dr$ with $(w-z)(s)=0$. Hence $\|e(t)\|\le(M_0\beta_0+(t-s)M_0\beta_1)\int_s^t\|e\|$, so $e\equiv0$: $u=w\in C^1$, $u'=v$.

\step{4}{Classical}
$u(t)=T(t-s)x+F(t)$ with $f=Bu$; both terms are in $D(A_0)$ and $A_0u=T(t-s)A_0x+F'-Bu=u'-Bu$.

\whyworks{differentiating the Duhamel formula once more shows that the mild solution is already the derivative of itself, by Gronwall.}
\end{proof}
\uses{none.}

\begin{dimcard}{B122}{VAR-Q18-CE}{A norm-continuous perturbation with no classical solution}{\pw}{ch:varcoef}
There are a Hilbert space $X$, $A_0$ skew-adjoint and $B\in C([0,T];\mathcal L(X))$ skew-adjoint at every $t$ such that for some $x_*\in D(A_0)$ and every $t$ with $\sin(nt)\ne0$, $U(t,0)x_*\notin D(A_0)$. So VAR-P11(e) is false without $C^1$. The sub-case of multiplication operators $B=-iV(t)\cdot$ with $V$ norm-continuous in $L^\infty$ stays open.
\end{dimcard}

\begin{proof}
\step{1}{Construction}
Let $A$ be skew-adjoint with unitary group $T(t)$, $x\in\mathcal H\setminus D(A)$, $n=\|x\|$, $X=\mathcal H\oplus\C$, $A_0=A\oplus0$, $f(t)=T(t)x$ and $B(t)(y,c)=(cf(t),-\langle y,f(t)\rangle)$.

\step{2}{Properties of $B$}
$\langle B(y,c),(y',c')\rangle=c\langle f,y'\rangle-\langle y,f\rangle\bar c'=-\langle(y,c),B(y',c')\rangle$: skew-adjoint. $\|B(t)-B(t')\|\le2\|T(t)x-T(t')x\|\to0$: continuous. It is not $C^1$, else $t\mapsto B(t)(0,1)=(T(t)x,0)$ would be differentiable, giving $x\in D(A)$. The datum is $x_*=(0,1)\in D(A_0)$.

\step{3}{Explicit solution}
Put $C(t)=\sin(nt)/n$, $y(t)=C(t)T(t)x$, $c(t)=C'(t)=\cos nt$. The first component of the mild equation is $\int_0^tT(t-r)c(r)f(r)dr=(\int_0^tc)T(t)x=y(t)$; the second is $c(t)=1-\int_0^t\langle y,f\rangle=1-n^2\int_0^tC$, which is the identity $\cos nt=1-n^2\int_0^t\sin(nr)/n\,dr$.

\step{4}{Conclusion}
By uniqueness $U(t,0)x_*=(C(t)T(t)x,\cos nt)$. For $0<t<\pi/n$, $C(t)\ne0$ and $T(t)x\notin D(A)$ since $D(A)$ is invariant under $T(\pm t)$.

\whyworks{the perturbation moves the datum along the very direction that is not in the domain, at the speed of a rank-$2$ rotation.}
\end{proof}
\uses{B121.}

\begin{dimcard}{B123}{VAR-Q10-EX}{Zeros of positive measure with empty interior}{\pw}{ch:varcoef}
There is $u\in C^\infty([0,1]\times\R)$ with $i\partial_tu+\partial_x^2u=0$ ($V=0$), $u(t,\cdot)$ entire in $x$ for each $t$, $\mathrm{int}\,Z(u)=\emptyset$ and $|Z(u)|>0$. So in the class $L^\infty$, $V=0$ and $V'=\mathbf 1_{Z(u)}$ are both laws ($V'u\equiv0$), while in the class $C$ the potential is identified ($\mathrm{int}\,Z=\emptyset$).
\end{dimcard}

\begin{proof}
\step{1}{A flat function}
Let $K\subset[0,1]$ be a fat Cantor set (closed, nowhere dense, $|K|>0$) with gaps $(\alpha_j,\alpha_j+\ell_j)$, $\sum\ell_j\le1$. Take $g_0(s)=e^{-s^{-2}}$ for $s>0$, $B(s)=g_0(s)g_0(1-s)$, with $|B^{(k)}|\le(2C_1)^k(k!)^{3/2}$, and $g=\sum_jg_j$, $g_j(t)=e^{-\ell_j^{-4}}B((t-\alpha_j)/\ell_j)$. Since $\sup_{\ell\in(0,1]}\ell^{-k}e^{-\ell^{-4}/2}\le C_2^k(k!)^{1/4}$ and $e^{-\ell^{-4}/2}\le\ell$, the series of derivatives converge uniformly and $|g^{(k)}|\le A^k(k!)^{7/4}$. Every $g_j^{(k)}$ vanishes on $K$, so $g^{(k)}=0$ on $K$, and $g>0$ on the gaps.

\step{2}{The Tychonoff series}
$u(t,x)=\sum_k(-i)^kg^{(k)}(t)x^{2k}/(2k)!$. As $(2k)!\ge(k!)^2$ the $k$-th term is $\le A^k|x|^{2k}/(k!)^{1/4}$, so the series converges for all $x$, as do those of the derivatives; $u$ is smooth, entire in $x$, and $i\partial_tu+\partial_x^2u=0$ termwise (telescoping).

\step{3}{Zeros}
For $t\in K$ all $g^{(k)}(t)=0$, so $u(t,\cdot)\equiv0$ and $Z(u)\supset K\times\R$, of positive measure. If an open $W\subset Z(u)$ existed it would contain a point with $t$ in a gap ($K$ is nowhere dense); there $u(t,\cdot)$ is entire and vanishes on an interval, hence $\equiv0$, against $u(t,0)=g(t)>0$.

\whyworks{Tychonoff's non-uniqueness for the heat-type equation gives a solution flat on a fat Cantor set but analytic in $x$.}
\end{proof}
\uses{none.}

\begin{dimcard}{B147}{S10COL-ii}{An inward vector field on a Lipschitz domain: $\bar\Omega\subset\Phi_\sigma(\Omega)$}{\pw}{ch:phase10}
Let $\Omega\subset\R^n$ be bounded and Lipschitz with charts $(C_i,R_i,g_i,L_i)$ covering $\partial\Omega$ and vertical directions $e_i$; let $\theta_i\in C_c^\infty(C_i)$, $\theta_i\ge0$, $\sum\theta_i>0$ on $\partial\Omega$, $W=\sum\theta_ie_i$, $V=-W$ with flow $\Phi_s$ and $M=\sup|V|$. Then for every $\sigma>0$: $\bar\Omega\subset\Phi_\sigma(\Omega)\subset\{\operatorname{dist}(\cdot,\bar\Omega)\le M\sigma\}$, and $\Phi_\sigma(\Omega)$ is open and diffeomorphic to $\Omega$ (same $b_1$). This is (COL-ii); it replaces the cited Grisvard/Hofmann--Mitrea--Taylor source in S10A-T3.
\end{dimcard}

\begin{proof}
\step{1}{The cone of inward directions}
For $p\in\partial\Omega$, $D(p)$ is the set of $w$ for which there are $r,a>0$ with $\operatorname{dist}(x+tu,\complement\Omega)\ge at$ for all $x\in\bar\Omega\cap B(p,r)$, $u\in B(w,r)$, $t\in(0,r]$. It is open and invariant under positive scaling (witness $(\rho,\lambda a)$, $\rho=\min(r,\lambda r,r/\lambda)$) and closed under sums: with witnesses $(r_k,a_k)$ for $w_k$, $\rho=\min(1,r_1,r_2,r_2/(2+|w_1|))$ and $u_k=w_k+\frac12(u-w_1-w_2)$, the point $x_1=x+tu_1$ lies in $\Omega$ and then $x_1+tu_2=x+tu$ has distance $\ge a_2t$. So $D(p)$ is a convex open cone, and $0\notin D(p)$.

\step{2}{Each chart cone is inward}
For a chart, $\Gamma=\{z:z_n>L|z'|\}\subset D(p)$ for $p\in\partial\Omega\cap C$. With $z\in\Gamma$, $\delta=z_n-L|z'|$, $d=\operatorname{dist}(p,\complement C)$, take $r\le1$ with $\sqrt{1+L^2}\,r\le\delta/2$ and $r(2+|z|)\le d/2$. For $x\in\bar\Omega\cap B(p,r)$ one has $x_n\ge g(x')$ (limit of points of $\Omega$), $y=x+tu\in C$, $u_n-L|u'|\ge\delta/2$ and $y_n-g(y')\ge t\delta/2$. Since $\operatorname{dist}(x,\complement\Omega)\ge\min(b/\sqrt{1+L^2},\operatorname{dist}(x,\complement C))$ for $b=x_n-g(x')$, one gets $\operatorname{dist}(y,\complement\Omega)\ge at$. In particular $e\in D(p)$.

\step{3}{The chart field is inward at every point}
$W(p)=\sum_{i\in I(p)}\theta_i(p)e_i$ with $e_i\in D(p)$ for $i\in I(p)$, so $W(p)\in D(p)$ by Step~1. No compatibility between overlapping charts is needed: $e_j$ may lie outside $\Gamma_i$ and both are in the intrinsic convex cone $D(p)$.

\step{4}{From a direction to trajectories}
If $W(p)\in D(p)$ with witness $(r,a)$, continuity of $W$ gives $\eta$ with $|W(y)-W(p)|\le r/2$ on $B(p,\eta)$; with $\rho_p=\min(r,\eta/(1+M))$, $\psi_t(x)=x+tu$, $u=\frac1t\int_0^tW(\psi_sx)\,ds\in B(W(p),r)$, so $\operatorname{dist}(\psi_t(x),\complement\Omega)\ge at$, where $\psi_t=\Phi_{-t}$.

\step{5}{Uniform inclusion}
Cover the compact $\partial\Omega$ by balls $B(p_j,\rho_{p_j}/2)$; let $\sigma_1=\min\rho_{p_j}/2$, $c=\min a_{p_j}$. Points of $\bar\Omega$ within $\sigma_1$ of $\partial\Omega$ are handled by Step~4; farther points have $\operatorname{dist}(x,\complement\Omega)\ge\sigma_1$ and move at most $Mt$. With $\sigma_0=\min(\sigma_1,\sigma_1/(M+c))$ one has $\operatorname{dist}(\Phi_{-\sigma}(x),\complement\Omega)\ge c\sigma$ for $\sigma\le\sigma_0$.

\step{6}{All $s>0$}
$\psi_t(\Omega)\subset\Omega$ for $t\ge0$: otherwise take the first exit time $\tau>0$, $y=\psi_\tau(x)\in\partial\Omega$; for small $h$, $z=\psi_{\tau-h}(x)\in\Omega$ and Step~5 gives $y=\psi_h(z)\in\Omega$, contradiction. Then for $s>\sigma_0$ compose with $\psi_{\sigma_0}(x)\in\Omega$. Finally $|\Phi_\sigma(x)-x|\le M\sigma$ and $\Phi_\sigma$ is a diffeomorphism of $\R^n$.

\whyworks{the inward cone at each boundary point is convex and open, so a positive combination of chart directions is inward; compactness of the boundary turns this local fact into a uniform margin.}
\end{proof}
\uses{none.}

\begin{caution}[title={Where it stops}]
For a non-Lipschitz domain $D(p)$ can be empty: in the disc with a slit, at a point $p$ of the slit a transversal direction crosses the slit at time $h/|w\cdot n|$ and the estimate $\ge at$ fails, a tangent one keeps $p$ on the slit. This is why the slit disc is outside S10A-T3.
\end{caution}

\begin{dimcard}{B148}{P2-Q1-NL-model}{On a simply connected slit disc, harmonic functions are holomorphic plus anti-holomorphic}{\pw}{ch:phase10}
Let $P=\Delta$ ($Q_1=\partial_{\bar z}$, $Q_2=\partial_z$), $m\ge1$, and $\Omega$ bounded, simply connected, satisfying a Poincar\'e inequality (for instance the slit disc $D\setminus[0,1)$, which is not Lipschitz). Then $\ker_{H^m}\Delta=\mathcal H_1+\mathcal H_2$ as sets, stronger than density.
\end{dimcard}

\begin{proof}
\step{1}{A holomorphic derivative}
For $h\in\ker_{H^m}\Delta$, $h_z$ is holomorphic.

\step{2}{A primitive}
$\Omega$ is simply connected, so $h_z$ has a primitive $f$, and $h-f$ is anti-holomorphic ($\partial_z(h-f)=0$).

\step{3}{Regularity}
$f'=h_z\in H^{m-1}$, and by Poincar\'e $\|f-f_\Omega\|_{L^2}\le C\|\nabla f\|_{L^2}$; hence $f\in H^m$ and also $h-f\in H^m$.

\step{4}{Poincar\'e on the slit disc}
Cover $D\setminus[0,1)$ by three convex sets (upper half-disc, lower half-disc, $D\cap\{\operatorname{Re}z<0\}$) with pairwise overlaps of positive measure, and chain the means.

\whyworks{the obstruction of the annulus is the flux, which vanishes on a simply connected set; the only thing needed from the geometry is Poincar\'e.}
\end{proof}
\uses{B119.}

\begin{caution}[title={What is left}]
General $P$ (elliptic factors with $\lambda\ne0$, hyperbolic or transport ones, where the translations of Step~2$'$ of S10A-T3 fail), $m=0$, and domains without Poincar\'e (``rooms and corridors'') are open; there is no route for general $P$. The exact codimension for general $P$ (P2-Q1-COD) lacks continuous ``period'' functionals in $H^m$.
\end{caution}


\section{Statements recorded without proof here}

\listcard{S10A-C2, C12-Q2 = OPF-Q2}{The question reduces to triples with at least two non-compact pairs $\Pi_i\Pi_j$, outside real transports (OPR2-T2); proved in B118.}{reduction \pw; the question \openq}{S10-OPEN-A \S4.}
\listcard{S10A-P1 (ring, (i))}{Ring structure; (ii) holds with the topological premise from Rudin, RCA 13.11 (cited) or an elementary argument.}{\pw}{S10-OPEN-A \S3; the annulus is B119.}
\listcard{S10COL-L1, L2, L3, TA, TB, TC = (COL-ii)}{(COL-ii) proved in B147; the step B of S10A-T3 is replaced by a chart field $X=-\sum\theta_ie_i$, $\theta_i\ge0$. Theorem B (a field with $V\cdot\nu\ge\kappa$ a.e. is inward at every point, by integrating along almost every segment) is not rewritten here.}{\pw}{S10-COL; (COL-i) from cited source: Grisvard Lemma 1.5.1.9 (verified), Hofmann--Mitrea--Taylor Cor.~2.13.}
\listcard{S10B-C1}{The extension of OPF5B-T1 (density) to non-convex regions is \textbf{false} (a hole with an interior is enough); proved in B146.}{\falsest{} (extension)}{S10-OPEN-B; F10-B2 prevails; see B119.}
\listcard{P2-Q1 (universal form)}{False for $P=\Delta$ on every Lipschitz region with $b_1\ge1$ (A53, B119).}{\falsest}{S10A-P1.}
\listcard{P2-Q1-NL}{Model case: $\Delta$, $m\ge1$, with $L^{1,2}\subset L^2$ proved in B148; general non-Lipschitz case open.}{model \pw; general \openq}{F10-TRIAGE.}
\listcard{C12-P7 (``closed but not dense'')}{Differential example on the annulus, $m\ge1$.}{\pw}{A53, B119.}
\listcard{S10B-P3}{Density of $(SB)$ proved in B126; ``open and dense'' after the fix P3-c (bumps also at minima of value $0$; counterexample $u_0'=-\cos x-\cos2x$) proved in the triage.}{\pw}{F10-B2, F10-TRIAGE.}

\section{Nonlinear and scattering closures}

\begin{dimcard}{B124}{S10B-T1}{Hirota $N$-solitons are reflectionless potentials}{\pw}{ch:phase10}
Let $0<\kappa_1<\dots<\kappa_N$, $c_j>0$, $\varepsilon_j=(-1)^{j-1}$, $k_j=2\kappa_j$, and
\[
f(x)=\sum_{S\subseteq\{1..N\}}\ \prod_{j\in S}e^{\eta_j}\prod_{\substack{i<j\\ i,j\in S}}A_{ij},\qquad A_{ij}=\Bigl(\tfrac{k_i-k_j}{k_i+k_j}\Bigr)^2,\quad e^{\eta_j}=c_jd_je^{k_jx},
\]
with $d_j=\prod_{l\ne j}\frac{\kappa_j+\kappa_l}{|\kappa_j-\kappa_l|}$ and $u=2(\log f)''$. Then $u\in\mathcal S(\R)$, $b\equiv0$, $a(k)=\prod_j\frac{k-i\kappa_j}{k+i\kappa_j}$, and the discrete spectrum of $L_u=-\partial^2-u$ is $\{-\kappa_1^2,\dots,-\kappa_N^2\}$. The same holds at every time with $c_j(t)=c_j(0)e^{-8\kappa_j^3t}$. That $u(t)$ solves KdV is Hirota's theorem, taken as part of the definition of the multi-soliton class.
\end{dimcard}

\begin{proof}
\step{1}{The Wronskian of the half-solutions}
Put $\varphi_j=e^{-\kappa_jx}+\varepsilon_jc_je^{\kappa_jx}$, so $-\varphi_j''=-\kappa_j^2\varphi_j$, and $W=W(\varphi_1,\dots,\varphi_N)$. Multilinearity in the columns, choosing in column $j$ the term $e^{s_j\kappa_jx}$ with $s_j=+1$ iff $j\in S$, gives
\[
W=V_0\,e^{-\sum\kappa_jx}\sum_S\prod_{j\in S}(\varepsilon_jc_j)\,\rho(S)\,e^{2\sum_{j\in S}\kappa_jx},\qquad \rho(S)=\frac{V(s_1\kappa_1,\dots,s_N\kappa_N)}{V_0},
\]
where $V$ is the Vandermonde determinant and $V_0=V(-\kappa_1,\dots,-\kappa_N)\ne0$.

\step{2}{Sign and modulus of $\rho(S)$}
Factor by factor, for $i<j$: the pairs outside $S$ give $1$, the pairs inside $S$ give $-1$, a mixed pair gives $\pm(\kappa_i+\kappa_j)/(\kappa_j-\kappa_i)$. The number of negative factors is $\binom{|S|}2+\#\{i<j:\ i\notin S,\ j\in S\}\equiv\sum_{j\in S}(j-1)$, so $\operatorname{sign}\rho(S)=\prod_{j\in S}\varepsilon_j$. For the modulus, $|\rho(S)|=\prod_{i\in S,l\notin S}\frac{\kappa_i+\kappa_l}{|\kappa_i-\kappa_l|}=\prod_{j\in S}d_j\prod_{i<j\in S}\bigl(\tfrac{\kappa_i-\kappa_j}{\kappa_i+\kappa_j}\bigr)^2$. Hence $W=V_0e^{-\sum\kappa_jx}f(x)$ with $f>0$, so $W$ never vanishes and $2(\log W)''=u$.

\step{3}{Intertwining without quoting Crum}
Let $Ag=W(\varphi_1,\dots,\varphi_N,g)/W$: a monic operator of order $N$ with $\ker A=\operatorname{span}\{\varphi_j\}$ and $a_{N-1}=-W'/W$. Divide: $A(-\partial^2)=QA+R$ with $\operatorname{ord}R<N$. Both $A(-\partial^2)$ and $QA$ kill each $\varphi_j$, so $R\varphi_j=0$ for $N$ independent functions, hence $R=0$. Comparing the coefficients of $\partial^{N+2},\partial^{N+1},\partial^N$ gives $q_2=-1$, $q_1=0$, $q_0=2a_{N-1}'=-2(\log W)''$. So $Q=-\partial^2-u=L_u$ and $L_uA=A(-\partial^2)$.

\step{4}{Jost solutions and scattering data}
For real $k\ne0$ the function $\psi=Ae^{ikx}$ solves $L_u\psi=k^2\psi$. As $x\to+\infty$ the terms $S=\{1..N\}$ dominate and $A\to\prod_j(\partial-\kappa_j)$, so $\psi e^{-ikx}\to\prod_j(ik-\kappa_j)$. As $x\to-\infty$ the terms $S=\emptyset$ dominate and $\psi e^{-ikx}\to\prod_j(ik+\kappa_j)$. Thus $\psi_+=\psi/\prod_j(ik-\kappa_j)$ is the Jost solution at $+\infty$ and at $-\infty$ it has \emph{no} $e^{-ikx}$ component. With the convention $\psi_+(x,k)=a(k)\psi_-(x,-k)-b(-k)\psi_-(x,k)$ this means $b(-k)=0$ for all real $k\ne0$, so $b\equiv0$, and $a(k)=\prod_j\frac{k-i\kappa_j}{k+i\kappa_j}$ on $\R\setminus0$.

\step{5}{Zeros of $a$ and the spectrum}
Let $R(k)=\prod_j\frac{k-i\kappa_j}{k+i\kappa_j}$, analytic in $\operatorname{Im}k>0$ and continuous to $\R$. Then $h=a-R$ is analytic in the upper half-plane, continuous to $\R\setminus0$ and zero there. By Schwarz reflection $h$ extends across $\R\setminus0$ and vanishes on an interval, so $h\equiv0$. The zeros of $a$ in the upper half-plane are exactly $i\kappa_j$, and the discrete spectrum is $\{-\kappa_j^2\}$.

\step{6}{Schwartz decay}
At each end, $\log f$ equals a linear exponent plus $\log(1+\text{exponentially small})$. Hence $u=2(\log f)''$ and all its derivatives are $O(e^{-2\kappa_1|x|})$.

\whyworks{the Wronskian of the half-solutions is a positive sum of exponentials, so the Darboux--Crum operator $A$ is smooth, and its asymptotic form at the two ends is a constant-coefficient operator that only multiplies the plane wave: nothing is reflected.}
\end{proof}
\uses{OPF5D-L3, OPF5D-L4, OPF5D-L9.}

\begin{caution}[title=Checks and convention]
SymPy/mpmath: Cauchy minors for $N\le4$ ($4/4$), sign and modulus of $\rho$ for $N\le6$ (exact, $12/12$), Wronskian relative error $\le1.8\cdot10^{-41}$, residual of $(-\partial^2-u)\psi=k^2\psi$ at $10^{-196}$, asymptotics at $x=\pm40$ of order $e^{-40}$. The class of multi-solitons is \emph{defined} by Hirota's formula, as in INT-P5/P6; with this convention INT-T1(B) has no external source.
\end{caution}

\begin{dimcard}{B125}{S10B-T2}{A prime differential ideal in one indeterminate is $[A]:H_A^\infty$}{\pw}{ch:phase10}
Let $F$ be a differential field of characteristic $0$ and $J\subsetneq F\{y\}$ a non-zero prime differential ideal. Let $A\in J$ have minimal rank (order $n$ first, then degree $d$ in $y_n$), with initial $I_A$, separant $S_A=\partial A/\partial y_n$ and $H_A=I_AS_A$. Then $J=[A]:H_A^\infty$.
\end{dimcard}

\begin{proof}
\step{1}{$A$ is irreducible}
If $A=BC$ with $B,C\notin F$, each factor has strictly smaller rank (order $<n$, or order $n$ and degree $<d$). By primality one of them lies in $J$, against minimality. A factor in $F\setminus0$ is impossible since $J$ is proper.

\step{2}{$H_A\notin J$}
$I_A\ne0$ has order $<n$ or lies in $F\setminus0$. $S_A\ne0$ because $d\ge1$ and $\operatorname{char}F=0$, and it has degree $d-1$ in $y_n$. Both have rank smaller than $A$, so they are not in $J$; by primality no power $H_A^k$ is in $J$.

\step{3}{$\supseteq$}
If $H_A^kf\in[A]\subseteq J$, then $f\in J$, since $J$ is prime and $H_A^k\notin J$.

\step{4}{$\subseteq$}
Let $f\in J$. Ritt's reduction (differential pseudo-division by the derivatives $A^{(i)}$, whose initial is $S_A$, then algebraic pseudo-division by $A$) gives $a,b\ge0$ and $R$ reduced with respect to $A$ such that $S_A^aI_A^bf\equiv R\pmod{[A]}$. Then $R\in J$ and has smaller rank than $A$ unless $R=0$; minimality forces $R=0$. So $H_A^{\max(a,b)}f\in[A]$.

\whyworks{minimality of the rank leaves no room for a reduced element in $J$, so reduction by $A$ must end at zero.}
\end{proof}
\uses{none.}

\begin{caution}[title=Application to a solution]
For $u$ real-analytic on an interval, $J_u=\{P:P[u]\equiv0\}$ is prime (the ring of analytic functions on a connected set is a domain) and differential. If $u$ solves $u_t=F[u]$ then $(D_FP)[u]=\partial_t(P[u])\equiv0$, so $J_u$ is $D_F$-stable, and the theorem applies: this removes the dependence of OPF5C-T5 on a numbered citation of Ritt or Kolchin.
\end{caution}

\begin{dimcard}{B126}{S10B-P3}{Morse and (SB) is dense}{\pw}{ch:phase10}
Let $r\ge3$. The class of $u_0\in C^r(\mathbb T)$ with $u_0'$ Morse and satisfying (SB) (at each birth time $t_b$ there is at most one degenerate minimiser $(x_b,a_b)$, and $M_{t_b}(x_b)=\{a_b\}$) is dense. Together with openness and the density of Morse functions it is open and dense, so ``generic'' in OPA-Q2 is no longer conditional.
\end{dimcard}

\begin{proof}
\step{1}{Candidates and distinct critical values}
Let $m_1,\dots,m_p$ be the local minima of $u_0'$ with negative critical value; the candidate times are $t_i=-1/u_0'(m_i)$ and $x_i=m_i+t_iu_0(m_i)$. Take $\psi_i(x)=(x-m_i)\chi(x-m_i)$ with $\chi$ even, $\chi=1$ on $[-\rho,\rho]$, supported in $(-2\rho,2\rho)$. Then $\psi_i'=1$ near $m_i$ and $\psi_i''$ lives in an annulus without critical points where $|u_0''|\ge c$. For small $\varepsilon$ the perturbation $u_0+\varepsilon\psi_i$ moves only the critical value $u_0'(m_i)$. Choosing the $\varepsilon_i$ small and generic makes the critical values pairwise distinct and non-zero; so at each $t_i$ there is only one candidate. The set of candidates is unchanged by all further steps.

\step{2}{Uniqueness of the minimiser at $t_i$}
If $m_i$ is not a minimiser of $\Phi_{x_i}$, or $M_{t_i}(x_i)=\{m_i\}$, do nothing. Otherwise $m_i$ is an isolated strict local minimum of $\Phi_{x_i}$; choose $r_i$ so that the annulus $C_i=\{r_i\le|a-m_i|\le2r_i\}$ contains neither critical points of $u_0'$ nor other minimisers, and $\Phi_{x_i}\ge\Phi_{x_i}(m_i)+\eta_i$ on $C_i$. With $\Theta_i\in C^\infty$, $0\le\Theta_i\le1$, $\Theta_i=0$ on $|a-m_i|\le r_i$ and $=1$ outside $|a-m_i|<2r_i$, set $\theta_i=\Theta_i'$ (periodic, zero mean, supported in $C_i$). Replacing $u_0$ by $u_0+\varepsilon\theta_i$ changes $U_0$ by $\varepsilon\Theta_i$, leaves the jet at $m_i$ (hence $t_i,x_i$) untouched, and raises $\Phi_{x_i}$ away from $m_i$. Then $M_{t_i}(x_i)=\{m_i\}$.

\step{3}{Earlier conditions are open}
Each condition SB$_j$ is open in $C^3$: the first alternative is a strict inequality between continuous quantities; in the second the margin outside a neighbourhood of $m_j$ is stable, and inside it $\Phi'_{x_j}$ is strictly increasing except at one point because $m_j$ remains a non-degenerate minimum of $u_0'$ with value $-1/t_j$. Choose each $\varepsilon_i$ smaller than half the margin of every SB$_j$, $j<i$.

\step{4}{Conclusion}
After $p$ steps $\tilde u_0=u_0+\sum\varepsilon_i\psi_i+\sum\varepsilon_i'\theta_i$ is Morse with (SB), and $\|\tilde u_0-u_0\|_{C^r}\le\sum|\varepsilon_i|\|\psi_i\|_{C^r}+\sum|\varepsilon_i'|\|\theta_i\|_{C^r}$ is as small as desired. Translates $m_i+2\pi n$ do not interfere: a minimiser satisfies $X_{t_i}(a')=x_i$ while $X(m_i+2\pi n)=x_i+2\pi n$.

\whyworks{the candidate times are finitely many and are not moved by the bumps; each bump kills the competitor at its own time and, being chosen last-and-small, cannot undo what was done before.}
\end{proof}
\uses{OPF5A-L4, OPF5A-11, OPF5A-12.}

\begin{caution}[title=The case of minima of value zero (triage fix P3-c)]
A local minimum of $u_0'$ with critical value $0$ would correspond to an infinite time. For ``open and dense'' one also adds bumps at such minima: for $u_0'=-\cos x-\cos2x$ the value is attained at degenerate points, and perturbing it away is what the triage proved.
\end{caution}

\begin{dimcard}{B127}{S10B-L4}{The uniform shock rate with several births}{\modulo{[F1]}}{ch:phase10}
Under Morse$+$(SB), $u_0\in C^5$: for every birth time $t_b$ there are $\tau_b,r_b>0$ such that, uniformly for $|t-t_b|\le\tau_b$,
\[
\int_{|x-x_b|<r_b}|u-v|\,dx=2\nu\log\frac1{\max(|t-t_b|,\sqrt\nu)}+O(\nu),\qquad \int_{|x-x_b|\ge r_b}|u-v|\,dx=O(\nu).
\]
Hence on $(0,T]$: $\|u-v\|_{L^1}=2\nu\log\frac1{\max(d(t),\sqrt\nu)}+O(\nu)$, with $d(t)=\min_b|t-t_b|$, and $\sup_t\|u-v\|_{L^1}=\nu|\log\nu|+O(\nu)$, modulo the Oleinik$\iff$Kru\v zkov equivalence [F1].
\end{dimcard}

\begin{proof}
\step{1}{Margin instead of global convexity}
Fix $z_0$ below the Taylor radius and the distance from $a_b$ to the other critical points of $u_0'$. By (SB), $M_{t_b}(x_b)=\{a_b\}$; upper semicontinuity of $(t,x)\mapsto M_t(x)$ gives $r_b,\tau_b,\eta_0>0$ such that for $|t-t_b|\le\tau_b$, $|x-x_b|<2r_b$ every minimiser has $|\alpha|<z_0/2$ ($\alpha=a-a_b$) and $\Phi_x\ge\min\Phi_x+\eta_0$ for $|\alpha|\ge z_0$. The region $|\alpha|\ge z_0$ weighs $O(e^{-\eta_0/2\nu})$ in the Hopf integral.

\step{2}{Local geometry}
With $\eta=1-t/t_b$ and $\kappa_3=u_0'''(a_b)>0$, $J_t=\frac t{t_b}J_{t_b}+\eta$, $J_{t_b}(a_b+\alpha)=\frac{t_b\kappa_3}2\alpha^2(1+O(\alpha))$, hence $c(\eta+\alpha^2)\le J_t\le C(\eta+\alpha^2)$ for $\eta\ge0$; $\Phi_x$ is convex in the window for $t\le t_b$. The estimates (a)--(c) of OPF5A-L2 follow with the tail replaced by the margin $\eta_0$.

\step{3}{Before the birth, $t\in[t_b-\tau_b,t_b]$}
Inside the window, Steps 1--4 of OPF5A-T1 repeat. Step 2 (``$dx=J_t\,da$ with $X_t$ bijective'') is false globally after the first birth; it is replaced by the local bijection $X_t|_{W_b}$, strictly increasing because $J_t>0$ on $W_b\setminus\{a_b\}$. Since $tu_0''=\partial_aJ_t$ has constant sign on each side of $a_b$, $\int|\partial_aJ_t|/J_t\,da=2\log\frac1{\max(\sqrt\nu,\eta)}+O(1)$. Outside, the minimisers are non-degenerate: on the compact set $\mathcal K_b$ one has $J\ge2j_0>0$, and the box argument of OPF5A-L3 gives $O(\nu)$.

\step{4}{After the birth, $t\in[t_b,t_b+\tau_b]$}
This is the delicate point. With $\tau=t-t_b$, $a=a_b+\sqrt\tau A$, $x=X_t(a_b)-t\tau^{3/2}\xi$, the normal form $\Psi_\tau(\xi,A)=\tau^{-2}[\Phi_x(a_b+\sqrt\tau A)-\Phi_x(a_b)]$ converges in $C^4$ on compacts to $\xi A+\frac{\kappa_3}{24}A^4-\frac{A^2}{2t_b^2}$. Lemma U uses only the Taylor expansion of $u_0$ at $a_b$ to order $5$, hence is local. Large $|A|$ inside the window is controlled by the $A^4$ growth, outside by $\eta_0$. In the intermediate region ($|\xi|>\Xi$, $|x-x_b|<r_b$) the convexity of 4.2 fails for $t>t_b$; one uses instead the local estimates $\psi(z)\ge\max(0,c_1(z^4+\alpha^2z^2)-C\tau z^2)$, $\psi(z)\le C(z^4+\alpha^2z^2+\tau z^2)$ for $|z|\le z_0$.

\step{5}{Gluing}
$T_1'\cap(0,T]=\{t_1<\dots<t_B\}$ is finite. Take $\tau_b$ below half the gaps between births: at most one window is active at any $t$. Away from the windows OPF5A-L3 applies.

\whyworks{near a birth only a bounded window of labels $a$ matters; the margin $\eta_0$ replaces global convexity there, and away from births the minimisers are non-degenerate.}
\end{proof}
\uses{OPF5A-L2, OPF5A-L3, OPF5A-T1, OPF5A-T2, S10B-P3.}

\begin{caution}[title=Checked on an example]
For $u_0=\sin x+0.45\sin(2x+1)$ (second birth at $t_2=1.8370$) the ratio $(\|u-v\|_{L^1}-\text{pred})/\nu$ is stable at $-1.086,-1.085,-1.087$ for $\nu=10^{-3},3\cdot10^{-4},10^{-4}$, and bounded at $|t-t_2|=0.02$. The point of risk is Step~4: uniformity of Lemma U with the margin in place of global convexity.
\end{caution}

\begin{dimcard}{B128}{S10B-T3}{Uniform-in-time $H^{n-1}$ bounds from the Gardner laws}{\pw}{ch:phase10}
Let $u\in C(\R;\mathcal S(\R))$ solve $u_t+u_{xxx}+6uu_x=0$. For each $n\ge1$, $\sup_{t\in\R}\|u(t)\|_{H^{n-1}}\le C_n(\|u(0)\|_{H^{n-1}})<\infty$. Global existence in $\mathcal S$ for Schwartz data is the hypothesis (WP), taken from the literature.
\end{dimcard}

\begin{proof}
\step{1}{Gardner's identity}
If $u=w+\varepsilon w_x-\varepsilon^2w^2$ then
\[
u_t+6uu_x+u_{xxx}=(1+\varepsilon\partial-2\varepsilon^2w)\bigl(w_t+6(w-\varepsilon^2w^2)w_x+w_{xxx}\bigr).
\]
The formal inverse series $w=\sum\varepsilon^nw_n[u]$, with $w_0=u$, $w_n=-\partial w_{n-1}+\sum_{i+j=n-2}w_iw_j$, solves Gardner's equation, which is in conservation form $w_t=-\partial(3w^2-2\varepsilon^2w^3+w_{xx})$. Order by order $D_Fw_n$ is a total derivative, so $H_n=\int w_n[u]\,dx$ is conserved for $u\in C(\R;\mathcal S)$ (checked symbolically for $n\le8$).

\step{2}{Weights}
With $\operatorname{wt}\partial=1$, $\operatorname{wt}u=2$, $\operatorname{wt}\varepsilon=-1$, $w_n$ is homogeneous of weight $n+2$: a monomial $\prod y_{j_i}$ ($y_j=\partial^ju$) satisfies $\sum(j_i+2)=n+2$.

\step{3}{The quadratic part}
The linear part of $w_n$ is $(-\partial)^nu$, a total derivative. The quadratic part obeys $Q_n=-\partial Q_{n-1}+\sum_{i+j=n-2}L_iL_j$, and modulo total derivatives
\[
\int Q_{2n}=\sum_{i=0}^{2n-2}\int\partial^iu\,\partial^{2n-2-i}u=(-1)^{n+1}\int(\partial^{n-1}u)^2.
\]

\step{4}{Lowering the order of higher-degree terms}
A monomial of degree $d\ge3$ in $w_{2n}$ carries $D=2n+2-2d\le2n-4$ derivatives. If $j_1>j_2+1$ an integration by parts moves one derivative off the largest-order factor, and $\sum j_i^2$ drops strictly; the process ends with $j_1\le(D+1)/2$, i.e.\ $j_1\le n-2$.

\step{5}{Gagliardo--Nirenberg}
Put $m=n-1$, $X=\|\partial^mu\|_{L^2}$, $M_0=\|u\|_{L^2}$ (conserved). Then $\|\partial^ju\|_{2}\le CM_0^{1-j/m}X^{j/m}$ and $\|\partial^ju\|_\infty\le CM_0^{1-\theta_j}X^{\theta_j}$, $\theta_j=\frac{j+1/2}m<1$, so for a reduced monomial
\[
\Bigl|\int\prod y_{j_i}\Bigr|\le\|y_{j_1}\|_2\|y_{j_2}\|_2\prod_{i\ge3}\|y_{j_i}\|_\infty\le C(M_0)X^\gamma,\qquad \gamma=\frac{2n+1-3d/2}{n-1}<2.
\]

\step{6}{Conclusion}
$X(t)^2=\pm H_{2n}+O(1+X(t)^\gamma)$ with $\gamma<2$ and $H_{2n}$ constant, hence $\sup_tX(t)<\infty$; no induction is needed.

\whyworks{each Gardner law has the square of the top derivative as its only quadratic part, and every other term is of strictly lower growth.}
\end{proof}
\uses{OP-F5D (Class (Cl)), (WP).}

\begin{dimcard}{B129}{S10B-L8p}{The graded bound extends to $K=2N$}{\pw}{ch:phase10}
Let $Z_k$ be the variety of order-$(2N-1)$ jets of $N$-solitons and $I_k$ its ideal, which is radical. The ideal of order-$2N$ jets is
\[
P_k^{(2N)}=I_k\,\C[y_0..y_{2N}]+\bigl(y_{2N}-\Phi(k;y_0,\dots,y_{2N-1})\bigr),
\]
and the initial form of the added generator, $y_{2N}-\Phi(0;y)$, is monic in the new variable $y_{2N}$. Hence the proof of OPF5C-L8 applies verbatim and the graded lower bound at $k=0$ transfers to every $k$ for $K=2N$.
\end{dimcard}

\begin{proof}
\step{1}{The flow gives $\Phi$}
By OPF5C-L1 the map $\Psi$ is a weighted-homogeneous polynomial automorphism, and Mumford's flow $V'=(V^2+f)/U-(y_0+\lambda)U$ (exact division) is polynomial and weighted. On $Z_k$ it expresses $y_{2N}=\Phi(k;y_0..y_{2N-1})$, homogeneous of weight $2N+2$ in $(k,y)$.

\step{2}{The ideal}
($\supseteq$) both kinds of generator vanish on $2N$-jets of $N$-solitons. ($\subseteq$) the set of such jets is the graph of $\Phi$ over the dense subset of order-$(2N-1)$ jets, so its Zariski closure is the graph $\Gamma$ of $\Phi|_{Z_k}$. The projection $\Gamma\to Z_k$ is an isomorphism, so $\C[y_0..y_{2N}]/(I_k,y_{2N}-\Phi)\cong\C[Z_k]$ is reduced and $(I_k,y_{2N}-\Phi)$ is the ideal of $\Gamma$.

\step{3}{Regular sequence of initial forms}
The initial form of $y_{2N}-\Phi$ (maximal $y$-weight component) is $y_{2N}-\Phi(0;y)$; it is monic in $y_{2N}$, a variable not occurring in the other generators, so it is a non-zero-divisor modulo them and the sequence of initial forms stays regular. This is the only property of the generators used in L8.

\whyworks{adding one monic equation in a fresh variable cannot spoil regularity of the initial forms.}
\end{proof}
\uses{OPF5C-L1, OPF5C-L8, OPF5C-T1.}

\begin{dimcard}{B130}{S10B-P10}{Finite-test criteria exist exactly when (I) fails or a finite hitting set exists}{\pw}{ch:phase10}
Let $\mathcal L=\{L_u:Pu=0,\ \Ann(u)\ne\mathfrak DP\}$ and call $\{T_1,\dots,T_s\}\subset\mathfrak D$ a finite-test criterion if, for every solution $u$ of $Pu=0$, $\Ann(u)=\mathfrak DP$ iff $T_iu\notin\mathcal T$ for all $i$. A finite-test criterion exists iff (I) \emph{no solution is identifying}, or $\mathcal L$ has a finite hitting set in $\mathfrak D\setminus\mathfrak DP$. In particular (I) is necessary for the theorem ``no finite test''.
\end{dimcard}

\begin{proof}
\step{1}{$\Leftarrow$, first case}
If (I) holds, take $T_1\in\mathfrak DP$: then $T_1u=0\in\mathcal T$ for all $u$, so the criterion declares every $u$ non-identifying, which is correct.

\step{2}{$\Leftarrow$, second case}
Let $T_1..T_s\notin\mathfrak DP$ hit every $L\in\mathcal L$. If $u$ is not identifying, some $T_i\in L_u$, i.e.\ $T_iu\in\mathcal T$. If $u$ is identifying, $L_u=\mathfrak DP$ and no $T_i\in\mathfrak DP$ gives $T_iu\notin\mathcal T$.

\step{3}{$\Rightarrow$}
Let $\{T_i\}$ be a correct criterion and suppose (I) fails, i.e.\ an identifying $u^*$ exists. If some $T_i\in\mathfrak DP$ the criterion fails on $u^*$. So all $T_i\notin\mathfrak DP$, and for each non-identifying $u$ some $T_iu\in\mathcal T$, i.e.\ $T_i\in L_u$: a finite hitting set.

\whyworks{a test either detects every non-identifying solution or it is useless; the detection set is exactly a hitting set of the annihilator ideals.}
\end{proof}
\uses{OP-ALG-T4, OPF5B-T6.}

\begin{caution}[title=What it does and does not give]
A condition such as ``infinitely many $u_\lambda$ with pairwise intersecting $L_{u_\lambda}$ in $\mathfrak DP$'' (any $T\notin\mathfrak DP$ is in at most one of them) is only \emph{sufficient} for the absence of a finite hitting set; realisability of the atoms of a repeated socle is not shown to be equivalent. OPR2-Q3 stays open in the exact form: with repeated socle, is (I) true and does $\mathcal L$ have no finite hitting set?
\end{caution}

\begin{dimcard}{B131}{S10B-P11}{Localised compact interaction closes the sum on a polygon}{\pw}{ch:phase10}
Let $\Omega\subset\R^2$ be a convex bounded polygon with vertices $v_1..v_V$, $Q_j=n_j\cdot\nabla+d_j$ with pairwise non-parallel $n_j$, $K_j=\ker_{L^2}Q_j$ with orthogonal projections $P_j$. Let $\chi_0,\dots,\chi_V$ be smooth with $\sum\chi_\alpha^2=1$ on $\bar\Omega$, $\chi_v\equiv1$ near $v$, $\chi_0\equiv0$ near all vertices. Assume (L-int): $P_i\chi_0^2P_j$ is compact for $i\ne j$; and (L-v): at each vertex, $\|\sum_j\chi_vf_j\|^2\ge c_v\sum_j\|\chi_vf_j\|^2-\langle\mathcal K_vf,f\rangle$ for some compact self-adjoint $\mathcal K_v$ on $\mathcal H=K_1\times\dots\times K_N$. Then $K_1+\dots+K_N$ is closed.
\end{dimcard}

\begin{proof}
\step{1}{Partition of the norm}
For $f=(f_j)\in\mathcal H$ and $S f=\sum f_j$, $\|Sf\|^2=\sum_\alpha\|\sum_j\chi_\alpha f_j\|^2$.

\step{2}{Interior term}
$\|\sum_j\chi_0f_j\|^2=\sum_j\|\chi_0f_j\|^2+\sum_{i\ne j}\langle f_i,P_i\chi_0^2P_jf_j\rangle=\sum_j\|\chi_0f_j\|^2-\langle\mathcal K_0f,f\rangle$ with $(\mathcal K_0f)_i=-\sum_{j\ne i}P_i\chi_0^2P_jf_j$, compact by (L-int) and self-adjoint since $(P_i\chi_0^2P_j)^*=P_j\chi_0^2P_i$.

\step{3}{Vertex terms and sum}
Use (L-v). With $c=\min(1,c_v)$ and $\mathcal K=\mathcal K_0+\sum_v\mathcal K_v$ one gets $\|Sf\|^2\ge c\|f\|^2_{\mathcal H}-\langle\mathcal Kf,f\rangle$.

\step{4}{Fredholm}
Hence $\mathcal G+\mathcal K$ with $\mathcal G=S^*S$ is coercive, thus invertible, and $\mathcal G$ is Fredholm. On $(\ker S)^\perp$, $\|Sf\|^2\ge c'\|f\|^2$, so $\operatorname{ran}S$ is closed.

\whyworks{a compact error in a coercive inequality does not change the closedness of the range: it can only add a finite-dimensional kernel.}
\end{proof}
\uses{OPF5B-P8, OPR2-L1.}

\begin{caution}[title=What is not proved]
The passage from the tangent-cone model (OPF5B-P8: a direct closed sum for homogeneous transports) to (L-v) on the polygon needs $\chi_vf_j$ to lie in the model kernel modulo compact errors; the commutator $(n_j\cdot\nabla\chi_v)f_j$ and the zeroth-order terms $d_j$ are not controlled. (L-int) is also not proved in this generality. So OPR2-Q2$'$ ($m=0$, polygon) stays open, reduced to (L-int)$\wedge$(L-v).
\end{caution}

\begin{dimcard}{B132}{S10B-E8}{No autonomous ODE in $x$ iff relation (R) fails (modulo S4ID-P8)}{\condl}{ch:phase10}
Assume S4ID-P8 (hence [GT-u$_m$] and (WP)): for a solution $u\in C(\R;\mathcal S)$ of KdV with $u(0)$ analytic and $r\not\equiv0$, every relation $G[u(t)]\equiv0$ has zero linear part and $\mathrm{in}(G)\in I_{PW}$, where $\mathrm{in}(G)$ is the homogeneous component of minimal degree and $I_{PW}$ the ideal of jets of $v''=cv$. Let (R) say: no $G\ne0$ with zero linear part and $\mathrm{in}(G)\in I_{PW}$ satisfies $G[u(t)]\equiv0$ for all $t$. Then INT-Q1 has the answer ``no ODE'' iff (R) holds.
\end{dimcard}

\begin{proof}
\step{1}{$\Rightarrow$}
If no relation exists, none exists in the class of (R).

\step{2}{$\Leftarrow$}
Let $G\ne0$ with $G[u(t)]\equiv0$. By S4ID-P8 its linear part vanishes and $\mathrm{in}(G)\in I_{PW}$; so $G$ is in the class of (R), contradicting (R).

\whyworks{S4ID-P8 already confines every possible relation to the class of (R).}
\end{proof}
\uses{S4ID-P8, OPF5C-T6.}

\begin{caution}[title={It is an equivalence, not a reduction}]
The case ``$G$ with non-zero linear part'' is excluded by S4ID-P8 itself, not by OPF5C-T6 (which for order $\ge4$ gives only exponential tails, compatible with $r\not\equiv0$). Deciding (R) needs the sub-leading terms of the radiative asymptotics ($t^{-1/2}\log$ corrections, the Deift--Zhou phase); Airy's cubic relation (OP-O5) lies in this class. INT-Q1 stays open.
\end{caution}

\begin{dimcard}{B133}{S10B-V9}{Values of $\delta_N(K)$: the computed table and what it proves}{computer-assisted \pw{} mod (Nov)}{ch:phase10}
Let $\delta_N(K)$ be the minimal degree of a homogeneous relation among the graded functions of OP-F5C at $k=0$ (the lower bound transferred to every $k$ by L8, L8$'$). Computed: $\delta_N(2N)=N+1$ for $N=1,\dots,5$; $\delta_5(8)=9$; and, in the triage with an independent mod-$p$ implementation on Adler--Moser jets, $\delta_5(10)=6$, $\delta_5(9)=7$ (reproduced) and $\delta_6(12)=7$ (new), the latter depending on the hypothesis that Adler--Moser jets lie in the closure of the principal stratum. Together with the upper bounds $\delta_N(K)\le N+1$ for $K\ge2N$ (OPR2-P4) and $\delta_N(2N-2)\le N+4$ (OPF5C-P7(b)), these give equalities for the values listed. The case $N=6$ in general is not executed.
\end{dimcard}

\begin{proof}
\step{1}{Reduction to a graded problem}
By L8 and L8$'$ the lower bound at $k=0$ holds for all $k\in U_0$. So the conjectures $\delta_N(2N)=N+1$, $\delta_N(2N-1)=N+2$ and $\delta_N(2N-2)=N+4$ are equivalent to the absence of homogeneous relations of degree $\le N$, $\le N+1$, $\le N+3$ respectively among the graded functions.

\step{2}{Computation}
\begin{center}\small
\begin{tabular}{@{}cccc@{}}
\toprule
$N$ & $K$ & minimal degree & time\\ \midrule
1 & 2 & 2 & 0.0 s\\ 2 & 4 & 3 & 0.1 s\\ 3 & 6 & 4 & 0.7 s\\ 4 & 8 & 5 & 21.5 s\\ 5 & 10 & 6 & 3231 s\\ \bottomrule
\end{tabular}
\end{center}
Each row gives $\delta_N(2N)(k)\ge N+1$ for every $k$; with OPR2-P4, equality.

\step{3}{$\delta_5(8)$}
The completed run at $N=5$, $K=8$ gives $\delta_5(8)\ge9$ at $k=0$; L8 transfers it, and OPF5C-P7(b) gives $\le9$.

\whyworks{the parameter $k$ disappears after L8: what remains is a finite Gr\"obner computation in a graded ring.}
\end{proof}
\uses{OPF5C-L8, B129, OPR2-P4, OPF5C-P7(b).}

\begin{caution}[title=Status of the computations]
$N\le4$ was reproduced by the audit; $N=5$ and $\delta_5(8)=9$ were computed once in the package (the first $(6,12)$ run was stopped after 3\,h\,28\,min), and then reproduced in the triage in 26--161\,s by the independent implementation (timings: $\delta_5(10)$ 26\,s, $\delta_5(9)$ 88\,s, $\delta_5(8)$ 161\,s, $\delta_6(12)$ 195\,s). The triage runs with Hirota jets for $N=5$ and with $(6,11)$ were not concluded; the value $\delta_6(11)=8$ is recorded in the ledger with a written proof not reproduced here. The proofs are computer-assisted modulo L8/L8$'$ and the hypothesis on Adler--Moser jets. The general question is OPF5C-Q1, open.
\end{caution}

\listcard{S4ID-P8, OPR2-P7}{Conditional on [GT-u$_m$] and (WP): the radiative asymptotics with derivatives of the KdV solution, from Grunert--Teschl (arXiv:0807.5041, Thm~5.4) with a uniform error for $x/t$ in compacts.}{\condl}{F10-GTU. The use made of it is in B132.}
\listcard{[GT-u$_0$]}{Boundedness of the H\"older constants of GT (5.20)--(5.21) and of the decoupling $\varepsilon$ for $k_0=\sqrt{-x/12t}$ in compacts. Every step of GT~\S5 is continuous in $k_0$, so this is plausible, but it needs \S5 rewritten with explicit constants; the earlier ``reduction'' (U) is the same statement, not a simpler one, and was withdrawn.}{\openq}{F10-GTU; blocked on sources.}
\listcard{[GT-u$_m$]}{Justification of the differentiation in GT Rem.~5.5 (Deift--Zhou, CMP 165, not read).}{\openq}{blocked on a source not available offline.}


\section{Systems, variable coefficients, encoder and dictionary closures}

\begin{dimcard}{B134}{VAR-Q17a}{Lateral observability for $V=V_0(x)+v(t)$}{\pw}{ch:varcoef}
Let $V=V_0(x)+v(t)$ with $V_0\in L^\infty(J;\R)$, $v\in L^\infty(0,T;\R)$, in the frame of VAR-L2-cor. If the mild solution $\psi$ vanishes a.e.\ on $(0,T)\times\omega$, $\omega\subset J$ open, then $\psi_0=0$.
\end{dimcard}

\begin{proof}
\step{1}{Remove the time-dependent part}
$\psi=e^{-i\int_0^tv}\tilde\psi$ where $\tilde\psi$ solves with $V_0$ only; the zero sets coincide, so one may take $V=V_0(x)$.

\step{2}{Mollify in time}
$H=H_0+V_0$ is self-adjoint on $D(H_0)$ and $\psi(t)=e^{-itH}\psi_0$. For $\varphi\in C_c^\infty(-\eta,\eta)$ put $\psi_\varphi(t)=\int\varphi(s)\psi(t+s)\,ds=e^{-itH}\check\varphi(H)\psi_0$. Since $\check\varphi$ decays faster than any power, $\check\varphi(H)\psi_0\in D(H^\infty)\subset D(H_0)$.

\step{3}{Apply VAR-L2-cor}
$\psi_\varphi$ vanishes on $(\eta,T-\eta)\times\omega$ and $V_0$ does not depend on $t$, hence is $C^1$ in $t$; VAR-L2-cor on the shifted interval gives $\check\varphi(H)\psi_0=0$.

\step{4}{Limit}
With $\varphi_n$ an approximate identity, $\check\varphi_n(H)\psi_0\to\psi_0$ strongly, so $\psi_0=0$.

\whyworks{a time-independent potential commutes with time mollification, which moves the datum into the domain where the classical unique-continuation theorem applies.}
\end{proof}
\uses{VAR-L2-cor, B121.}

\begin{caution}[title={Part (b): $V(t,x)$ stays open}]
For genuinely time-dependent $V$, $\psi_\varphi$ solves $i\partial_t\psi_\varphi=(H_0+V(t))\psi_\varphi+R_\varphi$ with a source $R_\varphi=\int\varphi(s)[V(t+s)-V(t)]\psi(t+s)ds$. It vanishes on $\omega$ where $\psi=0$, but the source lives outside the observation set, and VAR-L2 is a \emph{homogeneous} unique-continuation statement. What would suffice is a Carleman estimate in $H^{-1}_t$ or a mollification commuting with $V(t)$. The only regularity lemma proved is: for $h\in C^1_c(0,T)$ and mild $\psi$, $\int h(s)U(s,0)\psi_0ds\in D(H_0)$. VAR-Q17(b) is reformulated, not reduced.
\end{caution}

\begin{dimcard}{B135}{VAR-P12d-quad}{No quadratic differential logarithm for a switched oscillator}{\pw}{ch:varcoef}
Let $\tau_1=\tau_2=\pi/(4\omega_1)$, $T=\tau_1+\tau_2$, $H(t)=-\partial_x^2+\omega(t)^2x^2$ on $L^2(\R)$ with $\omega=\omega_1$ on $[0,\tau_1)$ and $\omega=\omega_2=5\omega_1$ on $[\tau_1,T)$. There is no self-adjoint $K'=\frac12(ap^2+b(xp+px)+cx^2)+dx+ep+f$ ($a,\dots,f\in\R$, $p=-i\partial_x$) with $e^{-iTK'}=U(T,0)$ up to a phase. The same holds for a $C^\infty$ regularisation of $\omega$ close in $L^1$.
\end{dimcard}

\begin{proof}
\step{1}{Monodromy}
For quadratic Hamiltonians $U^*(x,p)U=M(x,p)$ with $M\in SL_2(\R)$ solving $\dot x=2p$, $\dot p=-2\omega^2x$.

\step{2}{Constant segments}
A segment of length $\tau$ at constant $\omega$ contributes $\begin{pmatrix}\cos\theta&\sin\theta/\omega\\-\omega\sin\theta&\cos\theta\end{pmatrix}$, $\theta=2\omega\tau$. Here $\theta_1=\pi/2$ and $\theta_2=5\pi/2$.

\step{3}{Trace}
$\operatorname{tr}M=-(\omega_1/\omega_2+\omega_2/\omega_1)=-26/5<-2$.

\step{4}{An affine quadratic $K'$ would give an exponential}
For such $K'$, $e^{iTK'}(x,p)e^{-iTK'}=e^{T\Xi}(x,p)+v$ with $\Xi\in\mathfrak{sl}_2(\R)$, and equality of the propagators up to phase forces $v=0$, $e^{T\Xi}=M$.

\step{5}{Contradiction}
$\exp(\mathfrak{sl}_2(\R))\subset\{\operatorname{tr}>-2\}\cup\{-I\}$: elliptic $\Xi$ gives $\operatorname{tr}=2\cos$, parabolic $2$, hyperbolic $2\cosh$. The trace is continuous in $\omega\in L^1$, so the regularised problem fails too.

\whyworks{a real symplectic matrix with trace below $-2$ is not an exponential in $\mathfrak{sl}_2(\R)$.}
\end{proof}
\uses{none.}

\begin{caution}[title=What it does not say]
At period $2T$ one has $\operatorname{tr}M^2=(26/5)^2-2>2$ and positive eigenvalues, so a quadratic logarithm exists for $U(2T,0)$ (up to phase and metaplectic sign). The potential $\omega^2x^2$ is unbounded, outside the $L^\infty$ class. The register must \emph{not} contain ``(d) false'': for bounded $V(t,x)$ and for non-quadratic logarithms (d) is open; the sub-question is whether a metaplectic operator with monodromy trace $<-2$ can be $e^{-iTK'}$ with $K'$ non-quadratic differential.
\end{caution}

\begin{dimcard}{B136}{VAR-C3Q2-R}{A rank criterion for the annihilator, and data that do not identify}{\pw}{ch:varcoef}
Let $A_2=\C[t,x]\langle\partial_t,\partial_x\rangle$, $B=\C[t,x]\langle\partial_x\rangle$ (a left Ore domain with fraction field $D_B$), $L=\partial_t^m+\sum_{j<m}b_j\partial_t^j$ monic with $b_j\in B$, $\operatorname{rank}_BM=\dim_{D_B}(D_B\otimes_BM)$. (a) If $Lu=0$, then $\operatorname{Ann}_{A_2}(u)=A_2L\iff\operatorname{rank}_B(A_2u)=m$. (b) $u(t,x)=f(x+t)$ with $f(y)=\exp(e^y)$ is not D-finite, and $(\partial_t-\partial_x)u=0$, so $u$ does not identify $\square=\partial_t^2-\partial_x^2$. (a$'$) If $A_2u$ is holonomic no $L\ne0$ is identified.
\end{dimcard}

\begin{proof}
\step{1}{Division}
$A_2=A_2L\oplus\bigoplus_{j<m}B\partial_t^j$ (S8-VAR \S5(a)).

\step{2}{(a) The map $\pi$}
$\pi:B^m\to A_2u$, $(\beta_j)\mapsto\sum\beta_j\partial_t^ju$, is onto since $A_2Lu=0$. For $Q\in\operatorname{Ann}(u)$ write $Q=Q'L+\sum\beta_j\partial_t^j$ with $\sum\beta_j\partial_t^ju=0$: $Q\in A_2L\iff\beta=0$. So $\operatorname{Ann}(u)=A_2L\iff\ker\pi=0$. If $\ker\pi\ne0$ it is a non-zero submodule of the torsion-free $B^m$, of rank $\ge1$, hence $\operatorname{rank}A_2u<m$; if $\ker\pi=0$ the rank is $m$.

\step{3}{(b) The example}
$u$ is not D-finite: otherwise $f=u(0,\cdot)$ would be D-finite in $x$, and entire solutions of linear ODEs with polynomial coefficients have finite order of growth, while $\exp(e^y)$ has infinite order. (The elementary proof via Touchard polynomials replaces the Wiman--Valiron citation.) And $(\partial_t-\partial_x)u=0$ with $\partial_t-\partial_x\notin A_2\square$ by $\partial_t$-order, so $\operatorname{Ann}(u)\supsetneq A_2\square$: the criterion ``not D-finite $\Rightarrow$ identifies'', true for $m=1$, fails for $m\ge2$.

\step{4}{(a$'$) Holonomic data}
$A_2/A_2L$ has Gelfand--Kirillov dimension $3$ while a holonomic module over $A_2$ has dimension $2$ (Bernstein's inequality, Bj\"ork, cited), so $A_2u\ne A_2/A_2L$.

\whyworks{identifying $L$ means that $A_2u$ is as large as the free module on $m$ generators over $B$; any extra relation lowers the rank.}
\end{proof}
\uses{S8-VAR \S5.}

\begin{caution}[title=Open]
The non-monic case (VAR-Q13) stays open: without monicity the decomposition of Step~1 fails, and the equivalence ``$\operatorname{Ann}(u)\subset A_2\cap D_BL\iff\ker\pi=0$'' stated first is withdrawn. Part (a$'$) rests on Bernstein--Bj\"ork (cited), later replaced for the specific claim by the rank criterion.
\end{caution}

\begin{dimcard}{B137}{SYS-DIRAC-POL}{Euler product for the Dirac $1{+}1$ flow: formal yes, fixed time no}{\pw}{ch:systems}
Let $A(\xi)=X+mY$, $X=i\xi\sigma_z$, $Y=-i\sigma_x$, $\omega=\sqrt{\xi^2+m^2}$, so $e^{tA}=e^{g_1X}e^{g_2mY}e^{g_1X}$ with $g_1=a/\xi$, $g_2=-b/m$, $b=-\arcsin(m\sin\omega t/\omega)$, $2a=\arctan(\xi\tan\omega t/\omega)$. (a) $g_1,g_2\in\Q[\xi,m][[t]]$ and the identity holds in $\mathrm{Mat}_2(\Q[\xi,m][[t]])$. (b) For real $t\ne0$, $m\ne0$ there are no $k$ and real polynomials $p_j,q_j\in\R[\xi]$ with $\prod_{j=1}^ke^{p_jX}e^{q_jmY}=e^{tA(\xi)}$ for all $\xi\in\R$.
\end{dimcard}

\begin{proof}
\step{1}{(a) Formal polynomiality}
$\sin\omega t/\omega$ and $\tan\omega t/\omega$ are series in $t$ with coefficients polynomial in $\omega^2=\xi^2+m^2$; $\arctan(\xi\,\cdot)$ and $\arcsin(m\,\cdot)$ are odd, so $a/\xi$ and $b/m$ are series with polynomial coefficients. The identity is analytic for small real $t$, hence holds between formal series. Each truncation is a product of flows of differential operators ($\xi\leftrightarrow-i\partial_x$).

\step{2}{Lemma E}
Let $R_i\in\R[\xi]$ with $R_i-R_j$ non-constant for $i\ne j$, $d_j\in\C$. If $S=\sum d_je^{iR_j}\to0$ as $\xi\to+\infty$, then all $d_j=0$. Indeed $\frac1X\int_X^{2X}|S|^2\to0$ and expanding gives $\sum|d_j|^2+\sum_{i\ne j}d_i\bar d_j\frac1X\int_X^{2X}e^{i(R_i-R_j)}$; for non-constant $Q$, $Q'$ is eventually monotone and non-zero, and integrating by parts, $|\int_X^{2X}e^{iQ}|\le1/|Q'(X)|+1/|Q'(2X)|+|1/Q'(2X)-1/Q'(X)|=O(1)$, so the mixed terms tend to $0$.

\step{3}{(b) Expansion}
$e^{pX}=\operatorname{diag}(e^{ip\xi},e^{-ip\xi})$ and $e^{qmY}=\cos(qm)I-i\sin(qm)\sigma_x$ is a combination of $e^{\pm iqm}$. So $W_{11}(\xi)=\sum c_je^{iP_j(\xi)}$ with $P_j\in\R[\xi]$ (group equal $P_j$ up to constants).

\step{4}{Asymptotics of the target}
As $\xi\to+\infty$, $\omega=\xi+m^2/2\xi+O(\xi^{-3})$ and $(e^{tA})_{11}=\cos\omega t+i\xi\sin\omega t/\omega=e^{i\xi t}+\frac{im^2t}{2\xi}e^{i\xi t}+O(\xi^{-2})$.

\step{5}{Contradiction}
If $W=e^{tA}$, $G=W_{11}-e^{i\xi t}$ has the form of Lemma E and tends to $0$, so $G\equiv0$; but $G=\frac{im^2t}{2\xi}e^{i\xi t}+O(\xi^{-2})\ne0$ for large $\xi$.

\whyworks{finite words in polynomial times produce only oscillating exponentials with polynomial phases, and the $1/\xi$ correction of the exact flow cannot be produced by them.}
\end{proof}
\uses{SYS-E3 (S8-SYS).}

\begin{caution}[title=Scope]
The ZXZ decomposition is a product of \emph{local} flows only as a formal series in $t$. At fixed time the factors are non-local Fourier multipliers. Complex polynomial times are not covered: Lemma E does not treat $e^{iP}$ with complex $P$; an ordering by growth would be needed. Open.
\end{caution}

\begin{dimcard}{B138}{II-T2bp-CRIT}{The pole of a spectral projector, and why the Jordan form does not decide it}{\pw}{ch:systems}
Let $\chi_{A(\xi)}=Q\cdot B$ with $Q,B$ coprime over $\C(\xi)$, $E(\xi)=\ker Q(A)$, $F(\xi)=\ker B(A)$ complementary near $\xi_0$, $P(\xi)$ the projector onto $E$ along $F$. (a) $P$ has a pole at $\xi_0$ iff $E_0\cap F_0\ne0$, where $E_0=\lim E$, $F_0=\lim F$ in the Grassmannian; equivalently $\det[e(\xi_0)\ f(\xi_0)]=0$ for saturated local bases. (b) The Jordan form of $A(\xi_0)$ does not decide: $A_1$, $A_2$ below have the same Jordan form at $0$, $P_1$ has no pole and $P_2$ has one. (c) $A=\operatorname{diag}(\xi,-\xi)$ has constant projectors despite an eigenvalue crossing at $0$.
\end{dimcard}

\begin{proof}
\step{1}{Saturated bases}
Columns of rational bases of $E$ and $F$ can be divided by powers of $(\xi-\xi_0)$ (Smith form over the DVR $\C[\xi]_{(\xi-\xi_0)}$), giving $e(\xi),f(\xi)$ analytic at $\xi_0$ with full-rank values; then $E_0=\operatorname{im}e(\xi_0)$, $F_0=\operatorname{im}f(\xi_0)$ and $P=[e\ 0][e\ f]^{-1}$.

\step{2}{(a)}
If $\det[e\ f](\xi_0)\ne0$, $P$ is analytic. If it vanishes, let $0\ne v\in E_0\cap F_0$. A bounded $P$ would have a limit $P_0$ along a sequence with $P_0e(\xi_0)=e(\xi_0)$ and $P_0f(\xi_0)=0$, so $P_0v=v$ and $P_0v=0$: contradiction.

\step{3}{(b) Counterexample}
\[
A_1=\begin{pmatrix}\xi&0&0\\0&-\xi&1\\0&0&-2\xi\end{pmatrix},\qquad A_2=\begin{pmatrix}\xi&1&0\\0&-\xi&0\\0&0&-2\xi\end{pmatrix},
\]
with $Q=\tau-\xi$, $B=(\tau+\xi)(\tau+2\xi)$. Both have Jordan form $J_2\oplus J_1$ at $\xi_0=0$. For $A_1$, $E=\langle e_1\rangle$, $F=\langle e_2,e_3\rangle$ for all $\xi$, so $P=\operatorname{diag}(1,0,0)$. For $A_2$, $P_{12}=1/(2\xi)$ and $F\ni(1,-2\xi,0)$, so $F_0\ni e_1\in E_0$.

\step{4}{(c)}
For $\operatorname{diag}(\xi,-\xi)$ the eigenvectors are $e_1,e_2$ for every $\xi$: the projectors are constant.

\whyworks{the pole is a statement about the limit position of two subspaces, which depends on the first-order jet of $A$ at $\xi_0$, not on $A(\xi_0)$ alone.}
\end{proof}
\uses{none.}

\begin{caution}[title=Statements that are false]
``Crossing of eigenvalues implies a pole'' is false for a derogatory $A$ (case (c)); and ``the pole depends only on $A(\xi_0)$ and the factorisation'' (Q-S8SYS-4) is false by (b). Both counterexamples are verified symbolically (idempotency, commutation).
\end{caution}

\begin{dimcard}{B139}{SYS-MAXWELL}{An identifying Maxwell field exists on every connected $\Omega$}{\pw}{ch:systems}
Let $R=\C[\tau,\xi_1,\xi_2,\xi_3]$, $P\in R^{8\times6}$ the Maxwell matrix (rows $\tau E-\xi\times B$, $\tau B+\xi\times E$, $\xi\cdot E$, $\xi\cdot B$) and $M_P=R^{1\times6}/R^{1\times8}P$. Then $\operatorname{soc}_{\mathfrak m}(M_P)=0$ for every maximal ideal $\mathfrak m$; by SYS-EX an identifying field $(E,B)$ exists on every connected $\Omega$. More generally, $\operatorname{pd}M_P\le n-2$ implies $\operatorname{soc}=0$.
\end{dimcard}

\begin{proof}
\step{1}{Syzygies}
$S_1=(\xi,0,-\tau,0)$ and $S_2=(0,\xi,0,-\tau)$, in blocks $(3,3,1,1)$, satisfy $S_iP=0$; e.g.\ $\xi\cdot(\tau E-\xi\times B)-\tau\,\xi\cdot E=0$.

\step{2}{Rank and kernel}
A $6\times6$ minor is non-zero, so $\operatorname{rank}P=6$ and the left kernel over $\operatorname{Frac}R$ is generated by $S_1,S_2$. If $sP=0$ over $R$, then $s=f_1S_1+f_2S_2$; the supports of $S_1,S_2$ are disjoint, so $f_1\xi_i,f_1\tau\in R$; writing $f_1=a/b$ reduced, $b$ divides $\xi_1,\xi_2,\xi_3,\tau$, hence is a unit. So the kernel is free of rank $2$.

\step{3}{Resolution}
$0\to R^2\to R^8\to R^6\to M_P\to0$ is exact: $\operatorname{pd}M_P\le2$.

\step{4}{Socle}
Let $k=R/\mathfrak m$ and $N=R^{1\times8}P$. $R_{\mathfrak m}$ is regular of dimension $4$, so $\operatorname{Ext}^i(k,R)=0$ for $i<4$. From $0\to N\to R^6\to M_P\to0$: $\operatorname{Hom}(k,M_P)\hookrightarrow\operatorname{Ext}^1(k,N)$. From $0\to R^2\to R^8\to N\to0$: $\operatorname{Ext}^1(k,R^8)\to\operatorname{Ext}^1(k,N)\to\operatorname{Ext}^2(k,R^2)$, with both ends zero. Hence $\operatorname{Hom}(k,M_P)=0$.

\whyworks{a module of projective dimension at most $n-2$ over a regular ring of dimension $n$ has no socle: depth is too high.}
\end{proof}
\uses{SYS-EX (S8-SYS), injectivity of $C^\infty(\R^n)$ (cited).}

\begin{dimcard}{B140}{ENC-T2cp-P}{No finite jet certifies a principal annihilator}{\pw}{ch:encoder}
Let $S=\C[\tau,\sigma]$, $\operatorname{Ann}_q(u)=\{P:P(\partial_t,L_q)u=0\}$, $u\in C^\infty(\Omega)$ with $\operatorname{Ann}_q(u)=(P_*)$, $P_*\ne0$. For every $(t_1,x_1)$ and every $\eta$ such that $J^\circ\setminus[x_1-\eta,x_1+\eta]$ has non-empty interior there is $u'\in C^\infty(\Omega)$ with $u'=u$ on $I\times(x_1-\eta,x_1+\eta)$ and $\operatorname{Ann}_q(u')=(P_*(\tau-\alpha))\ne(P_*)$. So no finite jet, nor data on a strip around $x_1$, certifies $\operatorname{Ann}_q(u)=(P_*)$, even among principal annihilators.
\end{dimcard}

\begin{proof}
\step{1}{Choice of $\alpha$}
$P_*(\alpha,\sigma)\not\equiv0$ for all but finitely many $\alpha$; take such an $\alpha$, so that $(\tau-\alpha)\nmid P_*$.

\step{2}{The perturbation}
Take $0\not\equiv f\in C_c^\infty(J^\circ\setminus[x_1-\eta,x_1+\eta])$, $h=e^{\alpha t}f$, $u'=u+h$. Then $P(\partial_t,L_q)h=e^{\alpha t}P(\alpha,L_q)f$. If $p(\sigma)=P(\alpha,\sigma)\ne0$, then $p(L_q)f=0$ is a linear ODE with constant non-zero leading coefficient and $f$ has compact support, so $f\equiv0$ by Cauchy uniqueness: contradiction. Hence $\operatorname{Ann}(h)=(\tau-\alpha)$.

\step{3}{$\supseteq$}
Obvious, the operators of $S$ commute.

\step{4}{$\subseteq$}
Let $R\in\operatorname{Ann}(u')$, $g=Ru=-Rh$. Then $(\tau-\alpha)g=0$, so $(\tau-\alpha)R\in(P_*)$, and since $\tau-\alpha$ is irreducible and does not divide $P_*$, $P_*\mid R$. Also $P_*g=RP_*u=0$, so $P_*Rh=0$ and $(\tau-\alpha)\mid R$. Coprimality in the factorial ring $S$ gives $R\in(P_*(\tau-\alpha))$.

\whyworks{a compactly supported exponential in time is invisible to a jet at $x_1$ but adds exactly one linear factor to the annihilator.}
\end{proof}
\uses{none.}

\begin{dimcard}{B141}{ENC-Q4-S}{Non-vacuous separation: different constants are not equivalent after encoding}{\pw}{ch:encoder}
Take $\mathcal V=\{q_0+c:c\in\R\}$, $q_0\in C^1(J)$ real, $J=[-a,a]$, and $W_c$ the Volterra encoder of $q_0+c$ normalised at $x_0=0$ ($c_V=0$, $N_V=I$). Write $u\sim u'$ iff $u'=e^{i\theta(t)}u$. (a) Intersections are non-empty: if $u\in\mathcal S_{q_0}$ then $e^{-ict}u\in\mathcal S_{q_0+c}$ and $u\sim e^{-ict}u$. (S): for $u\sim u'\ne0$ solutions with $V=q_0+c\ne V'=q_0+c'$, $[W_Vu]\ne[W_{V'}u']$ (classes modulo phases).
\end{dimcard}

\begin{proof}
\step{1}{A good time interval}
Let $G=\{t:u(t,\cdot)\not\equiv0\}$: it is open and non-empty. Fix a component $G_0$; outside $G$ both $u,u'$ vanish and give no information.

\step{2}{The phase of $u'$ is linear on $G_0$}
Near $t_1\in G_0$ choose $x_1$ with $u(t_1,x_1)\ne0$; then $e^{i\theta_1}=u'/u$ at $x_1$ is $C^1$, and lifts to a $C^1$ $\theta_1$ on the interval $G_0$. Subtracting the two equations on $G_0\times J$: $0=(-\theta_1'-\delta)e^{i\theta_1}u$ with $\delta=c'-c\ne0$, so $\theta_1=\kappa-\delta t$.

\step{3}{Assume $W_{V'}u'=e^{i\theta}W_Vu$}
Put $w=W_Vu$, $w'=W_{V'}u'$: free classical solutions since $WL_q=L_0W$. As in Step~2, $\theta$ lifts to $C^1$ on $G_0$, and $0=iw'_t+w'_{xx}=-\theta'e^{i\theta}w$, so $\theta\equiv\theta_0$ on $G_0$.

\step{4}{An eigenvalue problem}
$W_{V'}$ acts only in $x$: $e^{i\theta_0}w=W_{V'}u'=e^{i(\kappa-\delta t)}Zw$, $Z=W_{V'}W_V^{-1}$. So $Zw(t)=\mu(t)w(t)$, $\mu=e^{i(\theta_0-\kappa+\delta t)}$.

\step{5}{$Z=I+\mathcal V$ with $\mathcal V$ quasinilpotent}
By ENC-1 $W_V=I+\int_{-x}^{x}H(x,s)\,ds$ with bounded kernel, and the composition of two such operators has the same form since $|r|\le|s|\le|x|$. Thus $\mathcal Vf(x)=\int_{-|x|}^{|x|}k(x,r)f(r)dr$. With $G_n(\rho)=\max(|\mathcal V^nf(\rho)|,|\mathcal V^nf(-\rho)|)$, $G_{n+1}\le2\|k\|\int_0^\rho G_n$, so $\|\mathcal V^n\|\le(2\|k\|a)^n/n!$ and $\sigma(\mathcal V)=\{0\}$.

\step{6}{Conclusion}
$Zw=\mu w$ with $w\ne0$ gives $\mu-1\in\sigma(\mathcal V)=\{0\}$, so $\mu\equiv1$ on $G_0$ and $\delta=0$: contradiction.

\whyworks{the encoder is the identity plus a Volterra operator, so it has only the eigenvalue $1$; a shift in the potential would need a different eigenvalue.}
\end{proof}
\uses{ENC-1, ENC-N1$'$(v).}

\begin{caution}[title=Literal DM3 reading]
For $V=V(x)$ the intersection $\mathcal S_V\cap\mathcal S_{V'}$ is $\{0\}$ (VAR-T1(b)): (S) is vacuous. For $V(t,x)$ non-trivial intersections exist (B123: $u\in\mathcal S_0\cap\mathcal S_{\mathbf 1_Z}$) but $V'=\mathbf 1_Z$ is not regular and time-independent encoders do not exist (ENC-P0(b)). ENC-Q4 in this reading stays open.
\end{caution}

\begin{dimcard}{B142}{ENC-Q10-D}{A non-local explicit Darboux encoder for a time-dependent potential}{\pw}{ch:encoder}
Let $u$ be a free solution ($iu_t+u_{xx}=0$) without zeros on $\Omega=I\times J$, $U=u_x/u$, $V_1=-2\partial_x^2\log u=-2U_x$, $\Lambda f=f_x-Uf=u\,\partial_x(f/u)$. (i) $\Lambda$ maps free solutions to solutions of $L_0\chi=V_1\chi$, and $\ker\Lambda=\C u$ on each slice. (ii) For $\chi=\Lambda\varphi$ with $\varphi$ free, put $\kappa(t)=(\chi_x+U\chi)(t,x_0)/u(t,x_0)$ and
\[
\psi=u\Bigl[\int_{x_0}^{x}\frac{\chi}{u}\,dx'+c(t)\Bigr],\qquad c'=i\kappa.
\]
Then $\psi$ is free, $\Lambda\psi=\chi$, and $\psi$ is unique modulo $\C u$. The encoder $\chi\mapsto\psi$ is non-local in $x$ and in $t$.
\end{dimcard}

\begin{proof}
\step{1}{(i) by hand}
From $iu_t=-u_{xx}$ and $u_{xx}/u=U_x+U^2$: $iU_t=-U_{xx}-2UU_x$. For $\chi=\varphi_x-U\varphi$ with $i\varphi_t=-\varphi_{xx}$: $i\chi_t=-\varphi_{xxx}-iU_t\varphi+U\varphi_{xx}$ and $\chi_{xx}=\varphi_{xxx}-U_{xx}\varphi-2U_x\varphi_x-U\varphi_{xx}$. Summing, $L_0\chi=(U_{xx}+2UU_x)\varphi-U_{xx}\varphi-2U_x\varphi_x=-2U_x\chi=V_1\chi$. $\Lambda f=0\iff\partial_x(f/u)=0$ and $J$ is connected.

\step{2}{The value of $\kappa$}
With $W=u\varphi_x-u_x\varphi=u\chi$ and $\varphi_t=i\varphi_{xx}$, $u_t=iu_{xx}$: $\partial_t(\varphi/u)=iW_x/u^2=i(u_x\chi+u\chi_x)/u^2=i(\chi_x+U\chi)/u$, which depends only on $\chi$ and $u$.

\step{3}{$\psi$ is free exactly when $c'=i\kappa$}
$\Lambda\psi=u\partial_x(\psi/u)=\chi$. Then $\psi-\varphi=a(t)u$, $a=c-(\varphi/u)(\cdot,x_0)$, and $L_0(au)=ia'u$; since $u\ne0$, $\psi$ is free iff $a'=0$ iff $c'=\partial_t(\varphi/u)(\cdot,x_0)=i\kappa$.

\step{4}{Uniqueness}
Two free $\tilde\psi,\psi$ with the same $\Lambda$ differ by $b(t)u$ with $ib'u=0$, so $b$ is constant.

\whyworks{Darboux transformations relate free solutions to solutions of a Schr\"odinger equation with a time-dependent potential $V_1$; inverting $\Lambda$ requires an integral and a time integration constant fixed by freeness.}
\end{proof}
\uses{none.}

\begin{caution}[title={A real, regular, time-dependent case and what stays open}]
The Galilean boost $u_G=(x-2kt+a)e^{i(kx-k^2t)}$ of the free solution $x+a$ gives $V_1=2/(x-2kt+a)^2$: real, time-dependent and regular on $I\times J$ for $J\subset(2k\sup I-a,\infty)$. It is trivial (a symmetry of a stationary case), so what remains open is the existence of real time-dependent Darboux examples \emph{not} reducible by symmetry to stationary ones, and whether history dependence is admitted in the class $\mathfrak E$. Also, $1/u$ does \emph{not} solve the $V_1$-model when $u_{xx}\ne0$ (residual $2u_{xx}/u^2$).
\end{caution}

\begin{dimcard}{B143}{DICT-Q3}{The pointwise and formal tests do not decide in $C^\infty$}{\pw}{ch:dictionary}
Let the law be $A_u(y)\theta(y)=-L_0u(y)$ (RANGO). Necessary: (N1) $-L_0u(y)\in\operatorname{im}A_u(y)$ for all $y$; (N2) a formal solution at every point. (b) These do not suffice in $C^\infty$: for Schr\"odinger with real $V$ and $u(t,x)=e^{-1/x^2}$ ($x>0$), $0$ ($x\le0$), (N1) and (N2) hold and no continuous $V$ exists. (a) For real-analytic $u$, (N2) at every point implies a $C^\infty$ solution (Malgrange, cited).
\end{dimcard}

\begin{proof}
\step{1}{(N1)}
For $x\le0$: $u=u_{xx}=0$, so $0\in\operatorname{im}A_u$. For $x>0$: $A_u=u\ne0$ is bijective.

\step{2}{(N2)}
At $x=0$ every series of $u$ vanishes ($u$ is flat), so $\theta=0$ is a formal solution; elsewhere the case is trivial.

\step{3}{No continuous $V$}
On $x>0$ the law is forced: $V=u_{xx}/u=4/x^6-6/x^4\to+\infty$ as $x\to0^+$.

\step{4}{(a), analytic case}
By the Tougeron--Malgrange theorem on modules generated by analytic functions (Malgrange 1966, ch.~VI, not read). For $K=q=1$ the quotient $-L_0u/T_1u$ is meromorphic and formally holomorphic at each point, hence analytic.

\whyworks{flatness hides the singularity of the law from every formal and pointwise test.}
\end{proof}
\uses{none.}

\begin{caution}[title=Status]
DICT-Q3(c) is reformulated, not reduced: a law exists iff the $C^\infty$ division problem is solvable; for general $C^\infty$ $u$ there is no criterion. Open.
\end{caution}

\begin{dimcard}{B144}{DICT-Q4-C}{Misfit certificates with non-local atoms}{\pw{} (the infinite-rank implication from Amitsur: \cf)}{ch:dictionary}
Let $D_0,\dots,D_N:C^\infty(\Omega)\to\mathcal D'(\Omega)$ be linearly independent and continuous. TFAE: (1) for every $M$, functionals $\ell_1..\ell_M$ and $u\in F_n$ there is $v\in U_n$ with $\ell_j(v)=\ell_j(u)$ (A(iii)); (2) for every closed finite-codimension subspace $Z=\bigcap\ker\ell_j$ there is $w\in Z$ with $D_0w,\dots,D_Nw$ independent. Moreover (2) holds if no non-trivial combination $\sum c_iD_i$ has finite rank, and (1) fails if one does.
\end{dimcard}

\begin{proof}
\step{1}{(2)$\Rightarrow$(1)}
The polynomial $s\mapsto\det(\ell_j'(D_i(u+sw)))$ has leading coefficient $\det(\ell_j'(D_iw))\ne0$; locality is used only to produce $w$.

\step{2}{(1)$\Rightarrow$(2)}
Take $u=0\in F_n$; a $v$ with the same data lies in $Z$ and in $U_n$.

\step{3}{Infinite rank implies (2)}
If every $w\in Z$ had dependent $D_iw$, the restrictions $D_i|_Z$ would be locally linearly dependent. By Amitsur's lemma (Trans.\ AMS 114, 1965; Bre\v sar--\v Semrl 1999, cited, not read) some $c\ne0$ has $\operatorname{rank}(\sum c_iD_i|_Z)\le N$, so $\operatorname{rank}\sum c_iD_i\le N+\operatorname{codim}Z<\infty$.

\step{4}{Finite rank implies failure}
If $D=\sum c_iD_i\ne0$ has finite rank $r$, $Z=\ker D$ is the intersection of $r$ continuous functionals and the family is dependent on $Z$.

\step{5}{Explicit example}
On $\Omega=(0,1)$, $D_if=\psi_i\int_0^1f$ with $\psi_i=1,x,x^2$: the single datum $\ell(u)=\int u=0$ certifies the fit, since every $v$ with $\ell(v)=0$ has $D_iv=0$.

\whyworks{a certificate fails exactly when a combination of atoms collapses onto finitely many functionals.}
\end{proof}
\uses{S8-DICT \S3.3.}

\begin{caution}[title=Consequence]
Non-zero differential operators have infinite rank, which recovers A(iii). The same holds for non-zero Fourier multipliers and Volterra convolutions $k*$ with $k\not\equiv0$ (infinite rank by Titchmarsh), so for memory and symbols A(iii) holds conditionally on Amitsur's lemma. In the triage the infinite-rank implication was replaced by DICT-LOCDEP.
\end{caution}

\begin{dimcard}{B145}{DM-delta-minimax}{The minimax rate $\varepsilon^{1/3}$}{\pw}{ch:dictionary}
Let $\mathcal U=\mathcal U(M,\delta_0,K)$ be the $u\in C^3(\bar K)$, $K=I\times J$, with $|u|\ge\delta_0$ and $\|\partial^\alpha u\|_{L^\infty}\le M$ for $|\alpha|\le3$, $V_u=(iu_t+u_{xx})/u$, $\tilde u$ with $\|\tilde u-u\|_\infty\le\varepsilon$. Assume $M>\delta_0$. The minimax risk satisfies $c\varepsilon^{1/3}\le R(\varepsilon)\le C\varepsilon^{1/3}$ on $K'\Subset K$ for $\varepsilon\le\varepsilon_0$.
\end{dimcard}

\begin{proof}
\step{1}{Upper bound}
Let $\hat V(\tilde u)=V_{u^*}$ for any $u^*\in\mathcal U$ with $\|u^*-\tilde u\|\le\varepsilon$ ($u$ itself qualifies). Then $e=u-u^*$ has $\|e\|\le2\varepsilon$ and $\|e_{tt}\|,\|e_{xxx}\|\le2M$, and VAR-P3(c) gives $\|V_u-V_{u^*}\|\le C\delta_0^{-1}(1+\|V_{u^*}\|)\bigl((2\varepsilon)^{1/3}(2M)^{2/3}+(2\varepsilon)^{1/2}(2M)^{1/2}+2\varepsilon\bigr)$ with $\|V_{u^*}\|\le2M/\delta_0$; for small $\varepsilon$ the $\varepsilon^{1/3}$ term dominates.

\step{2}{Lower bound, two points}
Take $u_1\equiv(M+\delta_0)/2$ (this is where $M>\delta_0$ enters) and $x_0$ interior to $K'$. Put $h=\varepsilon\varphi((x-x_0)/\varepsilon^{1/3})$ with $\varphi\in C_c^\infty$, $\varphi(0)=0$, $\varphi''(0)\ne0$, $\|\varphi\|_\infty\le1$. Then $\|h\|_\infty\le\varepsilon$, $\|\partial_x^3h\|=\|\varphi'''\|=O(1)$, and after rescaling $\varphi$, $u_2=u_1+h\in\mathcal U$.

\step{3}{The gap}
At $(t,x_0)$, $h=0$, so $V_{u_2}-V_{u_1}=h_{xx}/u_1=\varepsilon^{1/3}\varphi''(0)/u_1$. The datum $\tilde u=u_1+h/2$ is $\varepsilon$-close to both; every estimator errs by at least $\frac12|V_{u_2}-V_{u_1}|\ge c\varepsilon^{1/3}$ on one of them.

\whyworks{a bump of height $\varepsilon$ and width $\varepsilon^{1/3}$ is invisible in the data yet changes the second derivative by $\varepsilon^{1/3}$.}
\end{proof}
\uses{VAR-P3(c), A39.}

\begin{caution}[title=Scope]
This covers the declared model only (complex $V$, no zeros, deterministic noise). It does not touch real $V$, the case with zeros, or stochastic noise.
\end{caution}

\begin{dimcard}{B146}{S10B-C1}{Density fails around a hole: the rectangle functional}{\pw}{ch:phase10}
Let $\Omega=(0,3)^2\setminus[1,2]^2$ (bounded, connected, Lipschitz, non-convex), $P=\partial_x\partial_y$, $Q_1=\partial_x$, $Q_2=\partial_y$. For every $m\ge0$, $\ker_{H^m(\Omega)}Q_1+\ker_{H^m(\Omega)}Q_2$ is not dense in $\ker_{H^m(\Omega)}P$, not even in $L^2$. So the convexity in OPF5B-T1 cannot be dropped.
\end{dimcard}

\begin{proof}
\step{1}{The datum}
Take $\varphi\in C^\infty(\R)$, $\varphi=0$ on $(-\infty,1]$, $\varphi=1$ on $[2,\infty)$, and $u(x,y)=\varphi(x)\varphi(y)$. Then $u_{xy}=\varphi'(x)\varphi'(y)$ is supported in $[1,2]^2$, disjoint from $\Omega$: $Pu=0$ on $\Omega$ and $u\in H^m(\Omega)$ for all $m$.

\step{2}{The functional}
With $I_1=(0.4,0.6)$, $I_2=(2.4,2.6)$ and $F\in L^2(\Omega)$,
\[
\Lambda(F)=25^2\Bigl[\int_{I_2\times I_2}F-\int_{I_1\times I_2}F-\int_{I_2\times I_1}F+\int_{I_1\times I_1}F\Bigr].
\]
The four squares lie in $\Omega$ and $\Lambda$ is continuous on $L^2$.

\step{3}{$\Lambda$ kills both kernels}
The strips $(0,3)\times I_j$ are connected and lie in $\Omega$. If $\partial_xF=0$ there, $F(x,y)=g_j(y)$ on each strip, so $\int_{I_2\times I_j}F=\int_{I_1\times I_j}F$ and $\Lambda(F)=0$. Symmetrically for $\partial_yF=0$.

\step{4}{It does not kill $u$}
$u=1$ on $I_2\times I_2$ and $0$ on the other three squares, so $\Lambda(u)=1$. Since $\Lambda$ vanishes on the $L^2$-closure of $\ker Q_1+\ker Q_2$, $u$ is not in it.

\whyworks{$\Lambda$ is the monodromy of the form $u_xdx+u_ydy$ around the hole, and it is zero on every sum of a function of $x$ and a function of $y$.}
\end{proof}
\uses{OPF5B-T1.}

\begin{caution}[title=Sharper statement]
The counterexample works as soon as there is a coordinate rectangle with its four edge-strips in $\Omega$ and a closed coordinate rectangle with non-empty interior in $\R^2\setminus\Omega$ inside it. For a \emph{point} hole $\Omega=(0,3)^2\setminus\{p\}$ the claim holds only for $m=0$ ($u=H(x-p_1)H(y-p_2)$, $u_{xy}=\delta_p$); for $m\ge1$, $\delta_p\notin H^{-1}(\R^2)$ forces $u_{xy}=0$ on the square and density holds. For which non-convex $\Omega$ density holds is open.
\end{caution}


\listcard{VAR-P11(e)}{Autonomous proof by Volterra and Gronwall.}{\pw}{B121. VAR-L2-cor then \pw{} (it was from a cited source).}
\listcard{VAR-Q18-CE}{Rank-$2$ counterexample, also inside $L^2(J)$ with Neumann $H_0$; ``P11(e) with $B$ merely continuous'' is false.}{\falsest}{B122. VAR-Q18 (multiplications) \openq.}
\listcard{VAR-Q10-EX, VAR-T1(c)-Q10}{A solution with $V=0$, $|Z|>0$ and $\mathrm{int}\,Z=\emptyset$ (Tychonoff plus a Cantor-type set).}{\pw}{B123; closes the caution of A38.}
\listcard{ENC-Q4 (reading DM3), ENC-Q10 (real residue)}{Open; for ENC-Q10, the Galilean boost already gives $V_1=2/(x-2kt+a)$.}{\openq}{F10-C.}
\listcard{DICT-Q4-C, DICT-U3d-v, DM-$\delta$-minimax}{See B144, B145 and Chapter~\ref{ch:cIV3}.}{\pw}{F10-C.}
\listcard{P1-Q6 (Phase 10 form)}{Proved with constants in A52 (S10A-T1, C1); the literal statement is false (reaction--diffusion $u_t=Du_{xx}+ru$ has $E^{add}$ continuous from $H^1$).}{\pw}{F10-A; see also A52.}
\listcard{C3-Q1}{Proved as a classification by data model: composition of P4-T5+P4-L7, OP-ALG-T3, OPR2-C2, OP-ALG-T1 mod (AC), CONV-T2; no new mathematics.}{\pw (classification)}{F10-A.}

\section{Integration, examples and formalisation}

\listcard{EX10-a, b, c, d, e}{Textbook examples checked against the theory: free Dirac $1+1$ ($13/13$), harmonic oscillator ($19/19$), P\"oschl--Teller ($18/18$), fractional heat ($13/13$), Maxwell in vacuum ($17/17$). Total $80/80$ PASS; three new arguments proved.}{checked on examples}{S10-INT; audit EX10 (ex10\_b\_nd $9/9$, ex10\_e\_k $13/13$, ex10\_e\_h $10/10$).}
\listcard{F10-LEAN}{Formalisation in Lean~4 + Mathlib: $37$ entries, of which $27$ faithful, $7$ weaker and $3$ different from the ledger statement, none vacuous; all $85$ \texttt{\#print axioms} use only the standard axioms.}{verified}{S10-LEAN; independent review S-CHECK-LEAN. Mathlib has no Weyl algebra and no $C_0$-semigroups.}
\listcard{F10-G9}{Points delegated by Gate 9: ALG-TK (Trotter--Kato convergence only for relatively bounded $V$), and the renaming of the lower bound in VAR-P3 as $\kappa$ (not $\delta$).}{\cf; renaming}{GATE9.}
